
% ============================================================
% PART 01
% IDENTITY AND MANIFESTATION
% ============================================================
%
% Book:
%   Ontology of Viability
%
% Book Identifier:
%   VS.B01
%
% Part:
%   Identity and Manifestation
%
% Part Identifier:
%   VS.B01.P01
%
% File:
%   00_part_01.tex
%
% Purpose:
%   Main organizational file for Part 01.
%   Introduces the ontological scope of the part
%   and organizes its conceptual, formal, and
%   comparative investigations.
%
% ============================================================


% ------------------------------------------------------------
% PART DECLARATION
% ------------------------------------------------------------

\part{Identity and Manifestation}
\label{part:identity-and-manifestation}


% ------------------------------------------------------------
% PART METADATA
% ------------------------------------------------------------

\def\CurrentPartID{VS.B01.P01}
\def\CurrentPartTitle{Identity and Manifestation}


% ------------------------------------------------------------
% PART OPENING
% ------------------------------------------------------------

\begin{quote}
\textit{
To exist is to be identity-bearing. The inquiry into
identity therefore begins not with the separation of
being from identity, but with the examination of their
ontological inseparability and the different dimensions
through which identity can be understood.
}
\end{quote}


% ------------------------------------------------------------
% PART INTRODUCTION
% ------------------------------------------------------------

\section*{Introduction}
\addcontentsline{toc}{section}{Introduction}

The first part of Ontology of Viability establishes
the conceptual foundations for examining identity
and manifestation as fundamental dimensions of being.

The inquiry begins from a foundational ontological
commitment: existence and identity are inseparable.
To exist is to be identity-bearing. Being and identity
are not two successive conditions, nor does either
possess ontological priority over the other.

This commitment does not eliminate the conceptual
distinction between being and identity. Rather, it
establishes the basis for examining identity as a
distinct ontological concern without separating it
from the existence of the entity whose identity is
under consideration.

The purpose of this part is to clarify the conceptual
structure of identity and to establish the relations
that must be examined before a more comprehensive
account of manifestation and the dynamics of being
can be developed.

The investigation proceeds through three complementary
domains.

First, the conceptual exposition examines identity
through its relations to determination, distinction,
self-sameness, individuation, continuity, and
manifestation. These investigations establish the
conceptual vocabulary and distinctions required for
the subsequent development of the ontology.

Second, the formal reconstruction examines the
principles, arguments, and formal relations that
can be reconstructed from the conceptual foundations.
Its purpose is to make the logical structure of the
developed account explicit while distinguishing
established commitments from propositions that
require further justification.

Third, the comparative studies examine relevant
ontological approaches and conceptual traditions.
Their purpose is not to replace the internal
development of the Viability Philosophy System with
external classifications, but to clarify similarities,
differences, conceptual dependencies, and unresolved
questions.

The scope of this part must also be carefully
delimited.

The present investigation does not determine the
origin or source of an entity's identity, the
ontological level from which that identity arises,
or whether and how an entity's identity manifests
an identity belonging to a higher ontological level.

Such questions require a prior account of the
ontological domains and the status of light.
Their systematic examination belongs to the
investigation of the dynamics of being in Part 04.

Accordingly, Part 01 establishes the conceptual
groundwork for the ontology of identity and
manifestation without prematurely resolving
questions whose explanation depends on the
subsequent development of the ontological
architecture.


% ============================================================
% 01. CONCEPTUAL EXPOSITION
% ============================================================

\section*{Conceptual Exposition}
\addcontentsline{toc}{section}{Conceptual Exposition}

The conceptual exposition establishes the principal
distinctions and relations required for the
ontological examination of identity.

It develops the initial conceptual framework
through five chapters.


% ------------------------------------------------------------
% CHAPTER 01
% IDENTITY
% ------------------------------------------------------------


% ============================================================
% CHAPTER 01
% IDENTITY
% ============================================================
%
% Book:
%   Ontology of Viability
%
% Book Identifier:
%   VS.B01
%
% Part:
%   Identity and Manifestation
%
% Part Identifier:
%   VS.B01.P01
%
% Chapter:
%   Identity
%
% Chapter Identifier:
%   VS.B01.P01.C01
%
% File:
%   chapter_01.tex
%
% Purpose:
%   Main chapter file.
%   This file organizes and inputs the independent
%   conceptual sections of Chapter 01.
%
% ============================================================


% ------------------------------------------------------------
% CHAPTER DECLARATION
% ------------------------------------------------------------

\chapter{Identity}
\label{ch:identity}


% ------------------------------------------------------------
% CHAPTER METADATA
% ------------------------------------------------------------

\def\CurrentChapterID{VS.B01.P01.C01}
\def\CurrentChapterTitle{Identity}


% ------------------------------------------------------------
% CHAPTER OPENING
% ------------------------------------------------------------

\begin{quote}
\textit{
The ontological inquiry into identity concerns the
determinate being of what exists, rather than merely
the recognition, description, or classification of
what is encountered.
}
\end{quote}


% ------------------------------------------------------------
% CHAPTER INTRODUCTION
% ------------------------------------------------------------

\section*{Introduction}
\addcontentsline{toc}{section}{Introduction}

Identity constitutes one of the foundational problems
of ontology because the intelligibility of being cannot
be separated from the question of what makes an entity
the entity that it is.

Within the Viability Philosophy System, existence and
identity are not conceived as two successive
determinations in which existence first occurs and
identity is subsequently acquired, or in which identity
is established independently before existence becomes
actual.

Rather, being and having identity are inseparable
aspects of ontological actuality. To exist is to be
identity-bearing. This inseparability does not eliminate
the conceptual distinction between existence and
identity; it establishes that neither can be treated as
an ontologically prior stage in the constitution of
the other.

The present chapter establishes the conceptual
foundations required for examining identity within the
ontology of viability. Its task is to clarify the
ontological problem of identity and examine its
relations to determination, distinction, self-sameness,
individuation, continuity, and ontological relations.

The chapter does not attempt to provide a complete
account of the origin, ontological levels, or dynamic
constitution of identity. Such an account requires
prior clarification of the domains and levels of being
and the ontological status of light. These questions
belong to the broader architecture developed in the
subsequent parts of this book, particularly the
investigation of the dynamics of being.

Accordingly, the present inquiry distinguishes between
the inseparability of being and identity and the
further question of how the identity of a particular
entity is constituted within the broader ontological
architecture.

The objective is to establish a coherent conceptual
foundation without prematurely determining matters
whose explanation depends on the subsequent
development of the ontology.


% ------------------------------------------------------------
% SECTION 01
% THE ONTOLOGICAL PROBLEM OF IDENTITY
% ------------------------------------------------------------

\input{sections/section_01_ontological_problem_of_identity}


% ------------------------------------------------------------
% SECTION 02
% BEING AND IDENTITY
% ------------------------------------------------------------

\input{sections/section_02_being_and_identity}


% ------------------------------------------------------------
% SECTION 03
% IDENTITY AND DETERMINATION
% ------------------------------------------------------------

\input{sections/section_03_identity_and_determination}


% ------------------------------------------------------------
% SECTION 04
% IDENTITY AND DISTINCTION
% ------------------------------------------------------------

\input{sections/section_04_identity_and_distinction}


% ------------------------------------------------------------
% SECTION 05
% IDENTITY AND SELF-SAMENESS
% ------------------------------------------------------------

\input{sections/section_05_identity_and_self_sameness}


% ------------------------------------------------------------
% SECTION 06
% IDENTITY AND INDIVIDUATION
% ------------------------------------------------------------

\input{sections/section_06_identity_and_individuation}


% ------------------------------------------------------------
% SECTION 07
% IDENTITY AND CONTINUITY
% ------------------------------------------------------------

\input{sections/section_07_identity_and_continuity}


% ------------------------------------------------------------
% SECTION 08
% IDENTITY AND ONTOLOGICAL RELATIONS
% ------------------------------------------------------------

\input{sections/section_08_identity_and_ontological_relations}


% ------------------------------------------------------------
% SECTION 09
% SCOPE AND LIMITS OF IDENTITY
% ------------------------------------------------------------

\input{sections/section_09_scope_and_limits_of_identity}


% ------------------------------------------------------------
% SECTION 10
% SYNTHESIS OF IDENTITY
% ------------------------------------------------------------

\input{sections/section_10_synthesis_of_identity}


% ------------------------------------------------------------
% CHAPTER CONCLUSION
% ------------------------------------------------------------

\section*{Chapter Conclusion}
\addcontentsline{toc}{section}{Chapter Conclusion}

This chapter establishes the initial conceptual
framework for the ontological examination of identity
within the Viability Philosophy System.

The central ontological commitment developed here is
the inseparability of being and identity. Existence
does not precede identity as an independently
completed stage, nor does identity precede existence
as an independently established determination.

The conceptual examination of identity, however,
requires distinctions among its various dimensions.
Determination, distinction, self-sameness,
individuation, continuity, and ontological relations
represent different questions that must be examined
without reducing them to one another.

The conclusions of this chapter provide the conceptual
groundwork for the subsequent examination of
invariance, unity and plurality, manifestation, and
the relation between identity and manifestation.

The origin and ontological levels of identity remain
outside the scope of the present chapter. Their
systematic explanation requires the prior development
of the ontology of the realms of being and the
ontological account of light, followed by an
examination of the dynamics through which being and
identity are articulated.

The chapter therefore establishes a foundation rather
than claiming to have exhausted the ontological problem
of identity.


% ============================================================
% END OF CHAPTER 01
% ============================================================



% ------------------------------------------------------------
% CHAPTER 02
% INVARIANCE
% ------------------------------------------------------------


% ============================================================
% CHAPTER 02
% INVARIANCE
% Book I: Ontology of Viability
% ============================================================

\chapter{Invariance}
\label{ch:invariance}

% ------------------------------------------------------------
% CHAPTER METADATA
% ------------------------------------------------------------

\newcommand{\ChapterTwoTitle}{Invariance}
\newcommand{\ChapterTwoID}{VS.B01.P01.C02}

% ------------------------------------------------------------
% OPENING QUOTATION
% ------------------------------------------------------------

\begin{quote}
Invariance is not the absence of change, but the persistence
of a determinate aspect in relation to a specified domain
of change.
\end{quote}

% ------------------------------------------------------------
% CHAPTER INTRODUCTION
% ------------------------------------------------------------

\section*{Introduction}
\addcontentsline{toc}{section}{Introduction}

The preceding chapter established the foundational ontological
commitment that being and identity are inseparable. Existence
does not first occur as an identity-neutral condition to which
identity is subsequently added. To exist is to be identity-bearing,
and the distinction between being and identity is conceptual
rather than a separation between two independently constituted
ontological conditions.

Having established this starting point, the present chapter
examines invariance as a distinct but related ontological
problem.

The question of invariance emerges whenever an existent is
considered in relation to difference, transformation, or change.
An existent may undergo certain changes while some determinate
aspect of it remains unchanged in a specified respect. The
possibility of identifying such persistence raises a question
that cannot be answered merely by appealing to identity:
what does it mean for something to remain invariant in relation
to a domain of change?

For the purposes of the present first edition, invariance is
approached relationally. It is provisionally defined as a
relation between an existent, or a determinate aspect of an
existent, and a specified domain of changes or transformations,
such that the existent or aspect remains unchanged in a
particular respect under those changes or transformations.

This working definition contains three analytically distinguishable
components:

\begin{enumerate}
    \item The subject of invariance: an existent or a determinate
    aspect of an existent.

    \item The domain of change: the specified changes,
    transformations, or operations in relation to which
    invariance is considered.

    \item The respect of invariance: the particular aspect,
    determination, or relation that remains unchanged within
    the specified domain.
\end{enumerate}

These components are introduced to clarify the relational
structure of the concept. They should not be interpreted as
three independently existing ontological entities, nor should
their analytical distinction be mistaken for an ontological
separation.

The relational definition also prevents invariance from being
identified prematurely with absolute immutability. An existent
may change in certain respects without every aspect of it
necessarily changing. Conversely, the persistence of a
particular aspect does not establish that the existent as a
whole is immutable.

Invariance must therefore be distinguished from identity,
self-sameness, continuity, stability, regularity, symmetry,
and conservation. These concepts may intersect with the
problem of invariance, but their meanings and ontological
functions cannot be assumed to be identical.

The chapter does not attempt to establish the ultimate ground
or origin of invariance. It does not assume that invariance
must arise from an underlying substance, an immutable essence,
a fixed structure, or a higher ontological determination.
Such explanations would require arguments that exceed the
present scope.

The chapter also distinguishes ontological invariance from
its epistemic recognition. The fact that a knower identifies
a feature as invariant does not, by itself, establish that
the feature derives its existence from the act of recognition.
Nor does the inability to identify an invariant feature
establish that no such feature exists.

Knowledge and non-knowledge remain intertwined within the
situated condition of the knower. The limits of recognition
must therefore be acknowledged without converting those
limits into claims about the limits of being itself.

The purpose of this chapter is consequently limited but
foundational: to clarify the relational meaning of invariance,
to distinguish it from neighboring concepts, to examine its
possible forms and domains, and to identify the limits of
what can be established at this stage of the ontology of
viability.

The possibility that invariance may possess a deeper
ontological significance beyond its relational formulation
remains open. The present definition is a methodological
starting point for the first edition, not a final declaration
about the ultimate nature of invariance. This question may
be reconsidered in subsequent revisions after the broader
ontological architecture of the book has been developed.

% ------------------------------------------------------------
% SECTION 01
% ------------------------------------------------------------


% ============================================================
% SECTION 01
% THE ONTOLOGICAL PROBLEM OF INVARIANCE
% Chapter 02: Invariance
% Book I: Ontology of Viability
% ============================================================

\section{The Ontological Problem of Invariance}
\label{sec:ontological-problem-of-invariance}

% ------------------------------------------------------------
% 1.1. THE QUESTION OF INVARIANCE
% ------------------------------------------------------------

The problem of invariance emerges when an existent is
considered in relation to a domain of difference or change.
An existent may be subject to different conditions,
transformations, or alterations, while something about it
remains unchanged in a particular respect. This possibility
raises a question that cannot be adequately addressed by
examining identity alone: what does it mean for something
to remain invariant in relation to change?

The question is ontological because it concerns the
determinations of an existent and the manner in which
these determinations stand in relation to possible changes.
It is not merely a question about the ability of an observer
to recognize similarities between different states, nor is
it limited to the identification of recurring patterns.
The epistemic recognition of invariance is one dimension
of the problem, but it does not exhaust its ontological
significance.

The problem must therefore be approached without assuming
that everything which remains unchanged is necessarily
the identity of the existent, or that everything which
changes necessarily constitutes a transformation of its
identity. Both assumptions would prematurely collapse
distinct ontological questions into one another.

The first task is to clarify the meaning of invariance
itself and to identify the conditions under which an
invariance claim can be meaningfully formulated.

% ------------------------------------------------------------
% 1.2. THE RELATIONAL CHARACTER OF INVARIANCE
% ------------------------------------------------------------

For the purposes of this first edition, invariance is
provisionally understood as a relation between an existent,
or a determinate aspect of an existent, and a specified
domain of changes or transformations, such that the
existent or aspect remains unchanged in a particular
respect under those changes or transformations.

This definition is relational because invariance cannot
be adequately specified without identifying the domain
in relation to which something is said to remain unchanged.
A statement that something is invariant without specifying
what it is invariant in relation to remains incomplete.

The relational formulation requires three analytically
distinguishable elements:

\begin{enumerate}
    \item \textbf{The subject of invariance:} the existent
    or determinate aspect whose invariance is being
    considered.

    \item \textbf{The domain of change:} the specified
    transformations, alterations, or variations in relation
    to which invariance is examined.

    \item \textbf{The respect of invariance:} the particular
    determination, feature, or relation that remains
    unchanged within that domain.
\end{enumerate}

These elements define the structure of an invariance claim.
They do not establish that the elements themselves are
independent ontological entities. Their distinction is
analytical and serves to prevent conceptual ambiguity.

Consider, for example, an existent whose spatial position
changes while a particular structural relation among its
components remains unchanged. In this case, the position
does not remain invariant, but the specified structural
relation may do so. The invariance claim is therefore
restricted to a particular determination and a particular
domain of change.

This example illustrates why invariance should not be
understood as the absence of all alteration. It is possible
for some determinations to change while others remain
unchanged in relation to the same transformation.

The example does not, however, establish that structural
relations are the ultimate ground of identity or that
structural invariance is universal. It only demonstrates
the need to specify the respect in which invariance is
being asserted.

% ------------------------------------------------------------
% 1.3. INVARIANCE AND THE PROBLEM OF CHANGE
% ------------------------------------------------------------

The concept of invariance becomes meaningful through its
relation to a specified domain of change. This does not
mean that change is ontologically prior to invariance,
nor that invariance exists only as a consequence of an
observer's comparison between different states.

Rather, the relational definition identifies the terms
required to formulate a claim of invariance. The domain
of change establishes the respect in which the question
of remaining unchanged is being considered.

A distinction must therefore be maintained between
the occurrence of change and the determination of what
remains unchanged in relation to it.

Change may affect different aspects of an existent in
different ways. Some aspects may undergo alteration,
others may remain unchanged within the specified domain,
and the status of still others may remain undetermined.
It would be unjustified to assume that every change
affects every determination of an existent equally.

Likewise, it would be unjustified to infer from the
persistence of one determination that the existent as
a whole remains unchanged.

The problem of invariance consequently requires a
differentiated examination of the existent rather than
a general opposition between change and permanence.

% ------------------------------------------------------------
% 1.4. INVARIANCE IS NOT ABSOLUTE IMMUTABILITY
% ------------------------------------------------------------

Invariance must be distinguished from immutability.
Immutability, in its strongest sense, indicates the
absence of change in the subject under consideration.
Invariance, as provisionally defined here, does not
require such an absence.

An existent may undergo changes in certain respects
while remaining invariant in another respect. The
invariance of a particular determination therefore
does not establish the immutability of the whole existent.

Conversely, the absence of an observed change does not
by itself establish invariance. Such an absence may
result from the limited scope of observation, the
restricted duration of examination, or the fact that
the relevant domain of transformations has not been
specified.

The distinction is important because the concept of
invariance must not be transformed into an unsupported
claim about absolute permanence.

At this stage, the analysis does not determine whether
absolute invariance is possible, whether it has a
specific ontological status, or whether it applies
to any existent in an unrestricted sense. These
questions remain open.

% ------------------------------------------------------------
% 1.5. INVARIANCE AND THE INSEPARABILITY OF BEING
% AND IDENTITY
% ------------------------------------------------------------

The relational definition of invariance must remain
consistent with the foundational commitment established
in the preceding chapter: being and identity are
inseparable.

The introduction of invariance does not imply that
an existent first possesses an identity and subsequently
acquires invariant features as additional ontological
properties. Nor does it imply that invariance constitutes
a prior foundation from which identity emerges.

An existent is considered here as already identity-bearing.
The analysis of invariance concerns a further question:
how particular determinations of that existent stand
in relation to specified changes.

This distinction prevents invariance from being
mistaken for a substitute for identity. The fact that
a determination remains unchanged under a particular
transformation does not, by itself, explain why the
existent is the existent that it is.

Similarly, the fact that an existent undergoes change
does not, by itself, establish that its identity has
ceased to be.

The relation between identity and invariance therefore
requires further examination. The present section
establishes only that the two concepts address distinct
but connected problems.

% ------------------------------------------------------------
% 1.6. THE EPISTEMIC AND ONTOLOGICAL DIMENSIONS
% OF THE PROBLEM
% ------------------------------------------------------------

Claims about invariance are often formulated through
comparison. An observer examines different conditions
or states and identifies something that appears to
remain unchanged. This comparative activity belongs
to the epistemic dimension of the problem.

However, the epistemic act of identifying an invariant
must be distinguished from the ontological status of
what is identified.

The recognition of an invariant is conditioned by
the knower's position, available distinctions, methods
of examination, and limits of knowledge. A feature
may remain unidentified because the knower lacks
the conceptual or empirical resources necessary to
recognize it.

The absence of recognition therefore cannot be treated
as proof of the absence of invariance.

At the same time, an epistemic claim of invariance
must remain open to examination and correction.
What appears invariant within one limited domain
may prove variable when a broader or different domain
of transformations is considered.

Knowledge and non-knowledge are consequently intertwined
in the investigation of invariance. The knower necessarily
occupies a situated position from which some determinations
can be recognized while others remain unknown.

This limitation does not reduce being to knowledge.
It requires a distinction between what is ontologically
the case and what can be established by a particular
act of knowing.

% ------------------------------------------------------------
% 1.7. THE SCOPE OF THE INITIAL DEFINITION
% ------------------------------------------------------------

The relational definition adopted in this chapter
is deliberately limited.

It establishes that a meaningful claim of invariance
must specify:

\begin{itemize}
    \item the existent or determination under consideration;
    \item the domain of changes or transformations;
    \item the respect in which something remains unchanged.
\end{itemize}

It does not establish:

\begin{itemize}
    \item that invariance is the ultimate ground of identity;
    \item that every existent necessarily possesses an
    invariant feature;
    \item that every invariant feature is essential;
    \item that invariance is equivalent to permanence
    or immutability;
    \item that the ground of invariance must be a substance,
    essence, structure, or higher ontological determination;
    \item that invariance has the same form in every
    ontological domain.
\end{itemize}

These restrictions are not deficiencies in the definition.
They are methodological safeguards against introducing
metaphysical commitments before the relevant arguments
have been developed.

The possibility that invariance may possess a deeper
ontological significance beyond its relational
formulation remains open for subsequent investigation
and for future revisions of the philosophical system.

% ------------------------------------------------------------
% 1.8. THE ONTOLOGICAL SIGNIFICANCE OF THE PROBLEM
% ------------------------------------------------------------

The problem of invariance is significant because it
makes possible a more precise examination of the
relationship between an existent and the changes
in relation to which its determinations are considered.

Without the concept of invariance, the analysis may
remain trapped between two inadequate alternatives:
either the existent is assumed to remain entirely
unchanged, or every alteration is treated as if it
necessarily affects the existent in its entirety.

The relational approach makes room for a more
differentiated account. It permits the examination
of what changes, what remains unchanged, and what
cannot yet be determined within a specified domain.

This does not resolve the problem of persistence,
the origin of identity, or the ontological foundation
of continuity. It provides a conceptual framework
within which these questions can be distinguished
rather than prematurely answered.

In this respect, invariance contributes to the
development of the ontology of viability by refining
the analysis of how identity-bearing existents may
stand in relation to difference and transformation.

Its significance lies not in providing an immediate
explanation of persistence, but in clarifying the
conditions under which claims about remaining
unchanged can be formulated and critically examined.

% ------------------------------------------------------------
% 1.9. TRANSITION TO THE FOLLOWING SECTIONS
% ------------------------------------------------------------

Having established the ontological problem and the
provisional relational definition of invariance,
the analysis must now examine its relation to the
concepts developed in the preceding chapter.

The first and most immediate comparison concerns
identity. Since identity and being have already been
established as inseparable, it is necessary to
determine how invariance can be discussed without
being confused with identity or treated as its
ontological foundation.

The following section therefore examines the
relationship between invariance and identity,
identifying their conceptual differences and the
limits of what can be inferred from their relation.

% ============================================================
% END OF SECTION 01
% ============================================================


% ------------------------------------------------------------
% SECTION 02
% ------------------------------------------------------------


% ============================================================
% SECTION 02
% INVARIANCE AND IDENTITY
% Chapter 02: Invariance
% Book I: Ontology of Viability
% ============================================================

\section{Invariance and Identity}
\label{sec:invariance-and-identity}

% ------------------------------------------------------------
% 2.1. THE NEED TO DISTINGUISH IDENTITY FROM INVARIANCE
% ------------------------------------------------------------

The relationship between invariance and identity is
foundational to the conceptual development of this chapter.
Both concepts concern the determinations of an existent,
but they address different questions and must not be treated
as interchangeable.

The preceding chapter established that being and identity
are ontologically inseparable. To exist is to be
identity-bearing; identity is not an additional property
that is subsequently attached to an otherwise
identity-neutral existent.

The present chapter does not revise this commitment.
Instead, it examines how invariance can be understood
in relation to an existent whose identity is already
ontologically inseparable from its being.

Identity concerns the identity-bearing determination
of an existent as the existent that it is. Invariance,
according to the provisional relational definition
adopted here, concerns the relation between an existent
or one of its determinate aspects and a specified
domain of changes, in relation to which something
remains unchanged in a particular respect.

The distinction is therefore not between two
independently existing entities. It is between two
conceptual questions concerning the same existent.

The question of identity cannot be replaced by the
question of invariance, and the question of invariance
cannot be answered merely by repeating the proposition
that an existent is identical with itself.

% ------------------------------------------------------------
% 2.2. IDENTITY DOES NOT REQUIRE AN IMMEDIATE
% CLAIM OF INVARIANCE
% ------------------------------------------------------------

The inseparability of being and identity does not,
by itself, establish which determinations of an
existent remain invariant under a specified domain
of change.

To affirm that an existent is identity-bearing is
not yet to identify the aspects that remain unchanged
when that existent undergoes transformation.

For example, consider an existent whose spatial
position, external configuration, or relational
arrangement changes. The fact that these changes
occur does not automatically establish that its
identity has been lost. However, neither does the
continuation of its identity automatically identify
a particular feature that remains invariant.

These are distinct questions:

\begin{enumerate}
    \item Whether the existent is identity-bearing.
    \item Which of its determinations change under
    a specified transformation.
    \item Which of its determinations remain unchanged
    in relation to that transformation.
    \item Whether the identified invariance has any
    bearing on the persistence or continuity of its identity.
\end{enumerate}

The fourth question cannot be answered in advance.
It requires an argument concerning the relationship
between the particular invariant determination and
the identity of the existent.

Consequently, identity must not be reduced to a
collection of features that happen to remain unchanged
within a particular domain of observation or transformation.

% ------------------------------------------------------------
% 2.3. INVARIANCE DOES NOT CONSTITUTE IDENTITY
% ------------------------------------------------------------

A further distinction concerns the temptation to
treat invariance as the foundation or constitutive
source of identity.

If an existent preserves a particular feature under
a specified transformation, it may be tempting to
conclude that this feature is what makes the existent
the existent that it is.

Such a conclusion is not justified by the observation
of invariance alone.

The fact that a determination remains unchanged
under a particular domain of transformations establishes
a relational claim about that determination. It does
not establish that the determination is the ontological
ground of the existent's identity.

The distinction becomes especially important when
the invariant feature is structural, relational,
functional, or quantitative. The persistence of such
a feature may be relevant to the analysis of identity,
but its relevance does not automatically make it
the source of identity.

The ontology of viability therefore avoids identifying
an invariant determination with the identity of the
existent unless a separate argument establishes
that relationship.

The origin and constitution of identity remain
questions that require further ontological development.
They cannot be resolved merely through the relational
definition of invariance.

% ------------------------------------------------------------
% 2.4. THE SAME EXISTENT AND THE SAME DETERMINATION
% ------------------------------------------------------------

A distinction must also be maintained between the
identity of an existent and the persistence of a
particular determination belonging to it.

An existent may be considered under different
conditions or transformations. In such cases,
the question of whether the existent remains the
same must be distinguished from the question of
whether a particular determination remains unchanged.

These questions may have different answers.

A determination may remain invariant while other
determinations change. Conversely, a particular
determination may change while the existent's
identity remains an open question or continues
without that determination.

It is therefore necessary to distinguish at least
three analytical levels:

\begin{enumerate}
    \item The identity-bearing existence of the
    existent as such.

    \item The particular determinations through
    which the existent is considered.

    \item The invariance or variation of those
    determinations in relation to a specified
    domain of change.
\end{enumerate}

This distinction does not imply that the existent
can be ontologically separated from its determinations.
It only prevents the identity of the existent from
being equated with the unchanged persistence of
every particular determination.

The ontological relationship between the identity
of an existent and the determinations through which
it is constituted remains a subject for further
investigation.

% ------------------------------------------------------------
% 2.5. INVARIANCE AND THE FORMAL PRINCIPLE OF
% SELF-IDENTITY
% ------------------------------------------------------------

Formal self-identity is commonly expressed through
the proposition:

\[
A = A
\]

This proposition expresses the formal identity of
a term with itself. It does not, by itself, establish
that every determination associated with the referent
of that term remains unchanged under every possible
transformation.

The formal expression of self-identity and the
relational formulation of invariance therefore
operate at different conceptual levels.

Self-identity concerns the formal relation of a
term to itself. Invariance concerns whether a
specified determination remains unchanged under
a specified domain of transformations.

The distinction can be represented schematically.
Let \(A\) denote an existent and let \(T\) denote
a transformation within a specified domain.
If \(F\) represents a particular determination
of \(A\), an invariance claim may be expressed as:

\[
F(T(A)) = F(A)
\]

This expression is only a schematic representation.
It assumes that \(F\), \(T\), and the relevant
comparison are sufficiently defined. It does not
establish that the same mathematical formulation
applies to every ontological case.

The expression indicates that the determination
under consideration remains unchanged under the
specified transformation. It does not establish
that the entire existent is unchanged, nor that
\(F\) constitutes its identity.

Formal self-identity and invariance must therefore
remain conceptually distinct.

% ------------------------------------------------------------
% 2.6. THE POSSIBILITY OF MULTIPLE INVARIANT
% DETERMINATIONS
% ------------------------------------------------------------

An existent may be examined in relation to different
domains of change and different respects of
determination. Consequently, an invariance claim
concerning one aspect does not automatically
determine the status of every other aspect.

For example, a particular transformation may
preserve a structural relation while altering
a spatial configuration. Another transformation
may preserve a quantitative determination while
changing a relational arrangement.

These possibilities show that invariance claims
must be restricted to the specific determinations
and transformation domains under consideration.

The analysis does not yet establish whether
multiple invariant determinations are independent,
mutually dependent, hierarchically ordered, or
constitutively related.

Nor does the identification of multiple invariant
determinations establish that an existent possesses
multiple identities.

The question of how identity may relate to different
determinations across different ontological levels
remains unresolved. It must not be settled merely
by multiplying the number of invariant aspects
identified in an analysis.

This restriction is particularly important for
the subsequent development of the ontology of
viability, where the relationship between identity,
determination, and the different realms of being
will require a more comprehensive account.

% ------------------------------------------------------------
% 2.7. INVARIANCE AND THE LIMITS OF IDENTITY
% CLAIMS
% ------------------------------------------------------------

The distinction between identity and invariance
also imposes limits on what can be inferred from
the observation of an invariant feature.

First, the persistence of a feature under one
specified transformation does not establish its
persistence under every possible transformation.

Second, the persistence of a feature does not
establish that it is essential to the identity
of the existent.

Third, the alteration of a feature does not
establish that the existent has lost its identity.

Fourth, the absence of an identified invariant
feature does not establish that the existent
lacks identity or that no invariance is present.

These restrictions prevent the analysis from
confusing the limits of an invariance claim with
the ontological status of identity.

An identity-bearing existent cannot be treated
as ontologically dependent upon the observer's
ability to identify its invariant determinations.
The observer may recognize some aspects while
remaining unaware of others.

This epistemic limitation must be acknowledged
without converting it into an ontological
conclusion.

% ------------------------------------------------------------
% 2.8. THE RELATIONSHIP BETWEEN IDENTITY AND
% INVARIANCE REMAINS OPEN
% ------------------------------------------------------------

The analysis so far establishes a distinction
between identity and invariance, but it does not
establish that the two are unrelated.

An invariant determination may be relevant to
the way an existent is identified, distinguished,
or understood across different conditions.
It may also be relevant to the examination of
continuity and persistence.

However, the precise ontological significance
of such a determination cannot be inferred
from invariance alone.

Several questions remain open:

\begin{itemize}
    \item Can a particular invariant determination
    be constitutive of an existent's identity?

    \item Can an existent remain identity-bearing
    while all the determinations examined within
    a particular domain undergo change?

    \item Is there a necessary relationship between
    identity and any specific form of invariance?

    \item Can different forms of invariance have
    different ontological roles?

    \item How should the relationship between
    identity and invariance be understood across
    different ontological realms?
\end{itemize}

These questions require further argument. The
present chapter does not answer them by assuming
that invariance is either wholly external to
identity or necessarily constitutive of it.

The methodological commitment is to preserve
the distinction while keeping their possible
ontological relationship open.

% ------------------------------------------------------------
% 2.9. SYNTHESIS
% ------------------------------------------------------------

The relationship between invariance and identity
can now be stated in a preliminary form.

Identity concerns the identity-bearing existence
of an existent. Invariance concerns the relation
between that existent, or one of its determinate
aspects, and a specified domain of changes,
in relation to which something remains unchanged
in a particular respect.

The two concepts are therefore distinct but
potentially connected.

The distinction prevents three reductions:

\begin{enumerate}
    \item Identity must not be reduced to the
    persistence of invariant features.

    \item Invariance must not be treated as the
    ontological source of identity.

    \item The recognition of invariance must not
    be confused with the existence of identity.
\end{enumerate}

At the same time, the analysis does not deny
that invariant determinations may contribute
to a more comprehensive explanation of identity.
It only requires that such a contribution be
established through further argument.

The foundational commitment that being and
identity are inseparable remains intact. The
relational definition of invariance operates
within this commitment and does not replace it.

% ------------------------------------------------------------
% 2.10. TRANSITION
% ------------------------------------------------------------

Having distinguished invariance from identity,
the analysis must now examine the relationship
between invariance and difference.

The concept of invariance requires the specification
of a domain of change or transformation. Such
a domain necessarily raises questions about
difference: what varies, what remains unchanged,
and in what respect can these conditions be
distinguished?

The following section examines whether invariance
can be understood without a clear account of
difference, and whether the persistence of a
determination in relation to some differences
can coexist with variation in other respects.

% ============================================================
% END OF SECTION 02
% ============================================================


% ------------------------------------------------------------
% SECTION 03
% ------------------------------------------------------------


% ============================================================
% SECTION 03
% INVARIANCE AND DIFFERENCE
% Chapter 02: Invariance
% Book I: Ontology of Viability
% ============================================================

\section{Invariance and Difference}
\label{sec:invariance-and-difference}

% ------------------------------------------------------------
% 3.1. THE PROBLEM OF DIFFERENCE IN THE ANALYSIS
% OF INVARIANCE
% ------------------------------------------------------------

The concept of invariance cannot be adequately examined
without considering its relationship to difference.
An invariance claim identifies something that remains
unchanged in a particular respect within a specified
domain of changes or transformations. The specification
of such a domain necessarily raises the question of
what changes, what differs, and how the relevant
differences are distinguished from what remains
unchanged.

Difference is therefore indispensable to the analytical
formulation of invariance. However, this does not mean
that difference is ontologically prior to invariance,
or that invariance is merely an effect produced by
the comparison of different states.

The relationship must be approached without assuming
that every difference negates invariance, or that
every instance of invariance eliminates difference.
These assumptions would reduce two distinct concepts
to mutually exclusive conditions.

An existent may exhibit differences in certain
determinations while remaining invariant in another
respect. The task is to clarify the conditions under
which these two descriptions can be applied without
contradiction.

The present section therefore examines the different
ways in which difference may be related to invariance
and establishes the conceptual restrictions necessary
for a coherent account of both.

% ------------------------------------------------------------
% 3.2. DIFFERENCE AS A SPECIFIED DETERMINATION
% ------------------------------------------------------------

Difference must first be distinguished from the
undifferentiated notion of change.

To say that two conditions differ is to identify
a distinction between them with respect to some
specified determination. The determination may
concern position, quantity, configuration, relation,
quality, or another aspect relevant to the analysis.

Difference is therefore not adequately expressed
by the mere assertion that two conditions are not
identical. A meaningful analysis must clarify
the respect in which the difference is being
considered.

For example, an existent may occupy different
spatial positions under two conditions. The
positions differ, but this difference does not
necessarily establish that every other determination
of the existent also differs.

Likewise, two configurations may differ in their
arrangement while preserving a specified relation
among their components.

These examples indicate that difference is often
determination-specific. A difference concerning
one aspect does not automatically extend to all
other aspects.

This restriction is essential for the analysis
of invariance because a claim of invariance must
identify the particular determination that remains
unchanged despite the differences occurring
within the specified domain.

% ------------------------------------------------------------
% 3.3. DIFFERENCE DOES NOT NECESSARILY NEGATE
% INVARIANCE
% ------------------------------------------------------------

The existence of difference does not necessarily
exclude invariance.

If an existent undergoes a transformation in which
one determination changes while another remains
unchanged, the two descriptions refer to different
respects of the same case.

The difference and the invariance are not contradictory
because they do not concern the same determination
under the same conditions.

This can be expressed schematically. Let \(A\)
represent an existent, and let \(T\) represent
a specified transformation. Let \(F\) and \(G\)
represent two distinct determinations considered
in relation to that existent.

It is possible that:

\[
F(T(A)) \neq F(A)
\]

while:

\[
G(T(A)) = G(A)
\]

The first expression indicates variation in
determination \(F\), while the second indicates
invariance in determination \(G\), within the
specified transformation.

These expressions are schematic and presuppose
that the relevant determinations and transformation
are sufficiently defined. They do not establish
that every ontological difference can be represented
mathematically in this way.

Their purpose is to clarify that difference and
invariance may coexist when they concern distinct
respects of an existent.

Consequently, the presence of difference cannot
be treated as sufficient evidence against every
claim of invariance.

% ------------------------------------------------------------
% 3.4. DIFFERENCE WITHIN THE SAME DETERMINATION
% ------------------------------------------------------------

A more difficult problem arises when the same
determination is considered under different
conditions.

If a particular determination changes under
a specified transformation, it cannot simultaneously
be described as unchanged under that same
transformation and in the same respect.

The relevant distinction is between two different
claims:

\begin{enumerate}
    \item A determination changes in relation
    to a specified transformation.

    \item A determination remains unchanged
    in relation to that same transformation.
\end{enumerate}

If both claims refer to precisely the same
determination, transformation, and respect,
they are incompatible.

However, the incompatibility of these claims
does not establish that the existent as a whole
has changed in every respect or that its identity
has been lost.

It only establishes that the particular
determination under consideration is not
invariant within the specified domain.

The analysis must therefore avoid transferring
a conclusion about one determination to the
entire existent without an additional argument.

% ------------------------------------------------------------
% 3.5. DIFFERENCE BETWEEN EXISTENTS AND
% DIFFERENCE WITHIN AN EXISTENT
% ------------------------------------------------------------

The concept of difference operates in more than
one analytical context.

First, difference may concern the distinction
between two or more existents. In this case,
the analysis concerns how one existent differs
from another.

Second, difference may concern distinct conditions
or determinations associated with one existent.
In this case, the analysis concerns variation
within the consideration of a single existent.

These two contexts must not be confused.

The distinction between two existents is not
equivalent to the difference between two conditions
of one existent. Likewise, the variation of
a determination within one existent does not
automatically establish a difference between
two independent existents.

The relational definition of invariance must
therefore specify whether its domain concerns
transformations of one existent, comparisons
between different existents, or another form
of variation.

Without this specification, an invariance claim
may confuse the persistence of a determination
within one existent with the similarity or
difference between multiple existents.

The ontological status of relations between
different existents, and the extent to which
such relations contribute to their identity,
remain questions requiring further examination.

% ------------------------------------------------------------
% 3.6. DIFFERENCE, COMPARISON, AND THE
% POSITION OF THE KNOWER
% ------------------------------------------------------------

Difference is frequently recognized through
comparison. A knower considers two or more
conditions and identifies a determination
that differs between them.

This comparative activity is epistemic.
It depends upon the distinctions available
to the knower, the criteria of comparison,
and the domain within which the comparison
is conducted.

However, the epistemic recognition of a
difference must not be confused with the
ontological existence of the difference itself.

A difference may exist without being recognized
by a particular knower. Conversely, an apparent
difference may result from an inadequate
description, an inappropriate comparison,
or a mistaken identification of the relevant
determinations.

The same restriction applies to invariance.
The recognition of something as unchanged
depends upon the identification of the relevant
domain and respect of comparison. Yet this
epistemic dependence does not establish that
invariance is created by the act of recognition.

Knowledge and non-knowledge are intertwined
within the situated condition of the knower.
The recognition of one difference may coexist
with ignorance of other differences, just as
the identification of one invariant determination
may coexist with the inability to recognize
other invariant or variable aspects.

This limitation must not be interpreted as
a deficiency that can be eliminated simply
by assuming an unrestricted observer. The
knower is itself situated within being and
does not occupy a position outside the
ontological conditions being investigated.

% ------------------------------------------------------------
% 3.7. DIFFERENCE AND THE SCOPE OF AN
% INVARIANCE CLAIM
% ------------------------------------------------------------

An invariance claim is always restricted by
the domain within which it is formulated.

A determination may remain unchanged under
one set of transformations while varying
under another. Consequently, the identification
of invariance within a limited domain cannot
be generalized automatically to every possible
domain of change.

The relationship can be represented schematically.

Let \(D_1\) and \(D_2\) represent two distinct
domains of transformations. A determination
\(F\) may satisfy:

\[
F(T(A)) = F(A)
\qquad
\text{for } T \in D_1
\]

without satisfying the same condition for
transformations belonging to \(D_2\).

The expression does not establish that the
two domains are necessarily independent or
that every transformation can be assigned
to a clearly bounded mathematical set. It
illustrates the limited scope of an invariance
claim.

The same determination may therefore be
invariant relative to one domain and variable
relative to another.

This possibility prevents the concept of
invariance from being interpreted as an
unrestricted assertion of permanence.

It also establishes the need to identify
the domain of comparison before making
general claims about what remains unchanged.

% ------------------------------------------------------------
% 3.8. DIFFERENCE DOES NOT NECESSARILY
% IMPLY THE LOSS OF IDENTITY
% ------------------------------------------------------------

The distinction between difference and
invariance has direct implications for
the analysis of identity.

If an existent changes in one or more
determinations, it does not follow automatically
that its identity has ceased to be.

The previous chapter established that
existence and identity are inseparable.
The present analysis must preserve this
commitment without prematurely explaining
how identity relates to every possible
transformation.

A difference between two conditions of
an existent may be compatible with the
continuation of its identity. However,
the existence of such compatibility in
some cases does not establish a universal
rule governing every transformation.

The precise relationship between identity,
difference, and the persistence of an
existent remains dependent upon questions
that cannot be settled by the concept of
invariance alone.

In particular, the present section does
not determine which changes are compatible
with the persistence of identity, what
conditions may constitute a loss of identity,
or whether these questions have the same
answer across different ontological domains.

These issues must remain open until the
broader structure of the ontology has
been sufficiently developed.

% ------------------------------------------------------------
% 3.9. DIFFERENCE AS A CONDITION OF
% CONCEPTUAL DISTINCTION
% ------------------------------------------------------------

The analysis of invariance requires the
ability to distinguish what varies from
what remains unchanged. Difference therefore
has an indispensable analytical function.

Without identifying relevant differences,
it would be impossible to specify the
domain within which invariance is being
asserted.

However, this analytical function must
not be confused with an ontological claim
that difference produces invariance.

The identification of a distinction makes
an invariance claim conceptually articulable.
It does not necessarily establish the
ontological origin of the invariant
determination.

Likewise, the fact that a determination
is identified through comparison does
not mean that its existence depends upon
the comparison.

The analysis must consequently preserve
three distinctions:

\begin{enumerate}
    \item Difference as an ontological
    distinction or variation.

    \item Comparison as an epistemic
    operation through which differences
    may be recognized.

    \item Invariance as a relation in which
    a specified determination remains
    unchanged within a defined domain
    of change.
\end{enumerate}

These concepts are related, but none
can be reduced to another without
additional argument.

% ------------------------------------------------------------
% 3.10. SYNTHESIS AND TRANSITION
% ------------------------------------------------------------

The relationship between invariance and
difference can now be stated more precisely.

Difference identifies variation or
distinction with respect to a specified
determination. Invariance identifies
the persistence of a particular
determination in relation to a specified
domain of changes or transformations.

Difference does not necessarily negate
invariance because an existent may vary
in one respect while remaining unchanged
in another.

At the same time, the persistence of
one determination does not eliminate
the differences occurring in other
determinations.

The analysis therefore rejects two
reductions:

\begin{enumerate}
    \item The reduction of invariance
    to the absence of all difference.

    \item The reduction of difference
    to the necessary negation of
    every form of invariance.
\end{enumerate}

The relational definition adopted in
this chapter makes it possible to
examine both concepts within a specified
domain without treating them as
mutually exclusive.

Nevertheless, the present analysis
does not explain the ontological origin
of either difference or invariance.
Nor does it establish whether every
existent necessarily exhibits invariant
determinations.

These questions remain open for the
subsequent development of the ontology
of viability.

The next section examines the relation
between invariance and change. Although
difference and change are closely
connected, they are not identical.
A more precise analysis of their
relationship is necessary to determine
how invariance can be understood
in relation to transformation.

% ============================================================
% END OF SECTION 03
% ============================================================


% ------------------------------------------------------------
% SECTION 04
% ------------------------------------------------------------


% ============================================================
% SECTION 04
% INVARIANCE AND CHANGE
% Chapter 02: Invariance
% Book I: Ontology of Viability
% ============================================================

\section{Invariance and Change}
\label{sec:invariance-and-change}

% ------------------------------------------------------------
% 4.1. THE PROBLEM OF CHANGE IN THE ANALYSIS
% OF INVARIANCE
% ------------------------------------------------------------

The relationship between invariance and change constitutes
one of the central problems of the present chapter. The
relational definition of invariance adopted in this first
edition requires the specification of a domain of changes
or transformations in relation to which an existent, or
one of its determinate aspects, remains unchanged in a
particular respect.

Change is therefore directly relevant to the formulation
of invariance. However, their relationship must not be
understood through a simple opposition between two
mutually exclusive conditions.

If invariance were defined as the complete absence of
change, it would become indistinguishable from absolute
immutability. Such a definition would exclude the
possibility that an existent could undergo transformation
while some of its determinations remain unchanged.

Conversely, if every occurrence of change were understood
as the negation of all invariance, it would become
impossible to distinguish between the aspects of an
existent that vary and those that remain unchanged.

The analysis must therefore move beyond the general
opposition between permanence and change. It must
identify the respects in which change occurs and
the respects in which invariance may be asserted.

The purpose of this section is to clarify this relationship
without prematurely determining the ultimate ontological
ground of change, persistence, or invariance.

% ------------------------------------------------------------
% 4.2. CHANGE AS A DETERMINATE ALTERATION
% ------------------------------------------------------------

Change cannot be adequately understood as an
undifferentiated event in which an existent simply
becomes different.

A meaningful account of change requires the
identification of what changes, in what respect
it changes, and under which conditions the alteration
is considered.

An existent may undergo changes in position,
configuration, quantity, quality, relation, or
other determinations. These changes are not
necessarily equivalent, and they do not necessarily
affect every aspect of the existent in the same way.

For the purposes of this analysis, change is
considered as an alteration in one or more
determinations of an existent under specified
conditions.

This formulation is analytical. It does not
claim that every form of change can be reduced
to an alteration of properties, nor does it
assume that the distinction between an existent
and its determinations is an ontological
separation.

The distinction allows the analysis to examine
change without assuming that the existent as
a whole must be understood as either entirely
unchanged or entirely transformed.

% ------------------------------------------------------------
% 4.3. CHANGE DOES NOT NECESSARILY ELIMINATE
% INVARIANCE
% ------------------------------------------------------------

An existent may undergo change while some
determination remains unchanged in relation
to the same transformation.

Let \(A\) represent an existent and \(T\)
a specified transformation. Let \(F\) and
\(G\) represent two determinations considered
in relation to \(A\).

It is possible that:

\[
F(T(A)) \neq F(A)
\]

while:

\[
G(T(A)) = G(A)
\]

In this schematic representation, determination
\(F\) changes while determination \(G\) remains
unchanged under the specified transformation.

The expressions illustrate that change and
invariance may coexist when they concern
different determinations of the same existent.

This coexistence does not establish that
\(G\) is essential to the identity of \(A\),
nor does it demonstrate that the identity
of \(A\) necessarily persists through every
possible transformation.

It establishes only that a change in one
determination does not logically require
a change in every other determination.

The distinction is important because the
scope of a change must not be extended
beyond what the relevant analysis can
justify.

% ------------------------------------------------------------
% 4.4. THE DIFFERENCE BETWEEN LOCAL AND
% GENERAL CLAIMS OF CHANGE
% ------------------------------------------------------------

A claim about change may concern a particular
determination or the existent considered
more generally.

These are different levels of analysis.

A local claim identifies an alteration in
a specified determination. A general claim
concerns the extent to which the existent
as a whole has undergone transformation.

The first does not automatically establish
the second.

For example, a change in the spatial position
of an existent establishes a difference in
position. It does not establish that every
structural, relational, or other determination
has also changed.

Likewise, the preservation of a particular
relation does not establish that the existent
as a whole has remained unchanged.

The analysis must therefore distinguish
between:

\begin{enumerate}
    \item Change in a particular determination.
    \item Change in several determinations.
    \item Transformation of the overall
    configuration of an existent.
    \item A possible transformation affecting
    the identity of the existent itself.
\end{enumerate}

These distinctions identify different questions.
They do not establish a universal hierarchy
among the listed forms of change.

In particular, the fourth question cannot
be answered merely by accumulating observations
about changes in individual determinations.
It requires an account of the relationship
between those determinations and the identity
of the existent.

% ------------------------------------------------------------
% 4.5. INVARIANCE IS DOMAIN-DEPENDENT
% ------------------------------------------------------------

An invariance claim is restricted by the
domain of transformations in relation to
which it is formulated.

A determination may remain unchanged under
one set of transformations while changing
under another. Consequently, invariance
must not be interpreted as an unrestricted
claim that a determination remains unchanged
under every conceivable condition.

Let \(D_1\) and \(D_2\) represent two
specified domains of transformations.
A determination \(F\) may remain unchanged
under transformations belonging to \(D_1\),
while undergoing alteration under
transformations belonging to \(D_2\).

Schematically:

\[
F(T(A)) = F(A)
\qquad
\text{for } T \in D_1
\]

but there may be a transformation
\(T' \in D_2\) such that:

\[
F(T'(A)) \neq F(A)
\]

These expressions represent a restricted
invariance claim. They do not establish
that the domains are necessarily independent
or that every transformation can be assigned
to a formally defined mathematical set.

The important point is that invariance
and change must be examined relative to
specified domains.

Without such specification, the statement
that something is invariant remains
ambiguous, while the statement that
something changes may be too general
to support a precise ontological conclusion.

% ------------------------------------------------------------
% 4.6. CHANGE AND THE QUESTION OF
% CONTINUITY
% ------------------------------------------------------------

The relationship between change and
continuity must be distinguished from
the relationship between change and
invariance.

An existent may undergo a sequence of
changes, and the question may arise as
to whether it continues through those
changes.

However, the continuity of an existent
does not necessarily require that every
determination remain invariant.

Similarly, the invariance of a particular
determination does not by itself establish
the continuity of the existent as a whole.

These concepts address related but
different problems:

\begin{itemize}
    \item Change concerns alteration in
    one or more determinations.

    \item Invariance concerns the persistence
    of a specified determination relative
    to a defined domain of change.

    \item Continuity concerns whether and
    how an existent continues through
    different conditions or phases.
\end{itemize}

The relationship among these concepts
cannot be settled by treating any one
of them as a substitute for the others.

The ontological ground of continuity
and its relation to identity remain
open questions. They require a more
comprehensive account than the present
relational analysis of invariance can
provide.

% ------------------------------------------------------------
% 4.7. CHANGE DOES NOT NECESSARILY
% IMPLY LOSS OF IDENTITY
% ------------------------------------------------------------

The occurrence of change does not
automatically establish the loss of
identity.

The preceding chapter established
that being and identity are inseparable.
This commitment prevents the analysis
from treating existence as an initially
identity-neutral condition that may
simply lose identity whenever one of
its determinations changes.

At the same time, this foundational
commitment does not provide an immediate
criterion for deciding which transformations
are compatible with the persistence
of a particular existent and which,
if any, may involve a transformation
of identity.

The analysis must therefore distinguish
between the occurrence of change and
the ontological interpretation of that
change.

A change in position, quantity, configuration,
or relation may be identified without
having established what that alteration
means for the identity of the existent.

The question of identity through change
must remain open until the relationship
between identity, determination, and
the dynamics of being has been developed
more fully.

It would be premature to explain this
relationship by assuming an immutable
substrate, a permanent essence, or a
fixed collection of properties.

% ------------------------------------------------------------
% 4.8. CHANGE, OBSERVATION, AND
	EPISTEMIC LIMITATIONS
% ------------------------------------------------------------

The recognition of change frequently
depends upon the comparison of different
conditions or states.

An observer identifies an alteration
by distinguishing between what is
observed at one point or under one
condition and what is observed under
another.

This comparative activity belongs
to the epistemic dimension of the
analysis. It requires criteria for
identifying the relevant determinations
and distinguishing actual alteration
from differences in description,
measurement, or perspective.

However, the epistemic identification
of change must not be confused with
the ontological occurrence of change.

An existent may undergo an alteration
that remains unknown to a particular
observer. Likewise, an apparent change
may result from an inadequate comparison
or an error in the identification of
the relevant determination.

The same distinction applies to
invariance. The observer's inability
to identify a change does not establish
that no change has occurred. Nor does
the inability to identify an invariant
feature establish that no such feature
exists.

Knowledge and non-knowledge are
intertwined in the situated condition
of the knower. The analysis must
acknowledge this condition without
reducing the ontological status of
change or invariance to the limits
of observation.

% ------------------------------------------------------------
% 4.9. THE LIMITS OF THE RELATIONAL
% ACCOUNT OF CHANGE
% ------------------------------------------------------------

The relational definition of invariance
makes it possible to distinguish
between what changes and what remains
unchanged within a specified domain.

Nevertheless, it does not explain
why change occurs, what ultimately
constitutes transformation, or how
the identity-bearing existence of
an existent is related to its changing
determinations.

Several questions therefore remain
outside the scope of the present
section:

\begin{itemize}
    \item What is the ontological ground
    of change?

    \item Are all forms of change
    reducible to alterations in
    determinations, or are there
    other ontological forms of
    transformation?

    \item How should the relation
    between identity and change
    be understood?

    \item Under what conditions,
    if any, does a transformation
    involve a change in identity?

    \item Is invariance merely relational,
    or does it disclose a deeper
    dimension of ontological
    determination?

    \item How do the different realms
    of being participate in or
    condition the possibilities
    of change and invariance?
\end{itemize}

These questions are not answered
by assuming that change is always
external to an existent or that
invariance is always internal to it.
Such distinctions would require
additional ontological justification.

The present analysis is limited to
clarifying the conceptual relationship
between change and invariance.

% ------------------------------------------------------------
% 4.10. SYNTHESIS AND TRANSITION
% ------------------------------------------------------------

The relationship between invariance
and change cannot be adequately
described as a simple opposition
between permanence and alteration.

Change concerns the alteration
of one or more determinations under
specified conditions. Invariance
concerns the relation between an
existent, or one of its determinate
aspects, and a specified domain of
transformations, within which
something remains unchanged in
a particular respect.

These concepts may therefore coexist
within the analysis of the same
existent, provided that the relevant
determinations and respects are
properly distinguished.

The present section establishes
the following restrictions:

\begin{enumerate}
    \item Change in one determination
    does not necessarily imply
    change in every determination.

    \item Invariance in one respect
    does not establish the immutability
    of the existent as a whole.

    \item Change does not automatically
    imply the loss of identity.

    \item The persistence of an invariant
    determination does not by itself
    establish the continuity of the
    existent.

    \item Every claim of invariance
    must specify the domain of
    transformations in relation
    to which it is asserted.
\end{enumerate}

The analysis has therefore clarified
the conceptual conditions under
which change and invariance may
be considered together without
reducing either concept to the other.

The next section examines the
relationship between invariance
and continuity. This examination
is necessary because the persistence
of a determination under change
and the continuation of an existent
through different conditions are
related problems, but they cannot
be assumed to have the same
ontological meaning.

% ============================================================
% END OF SECTION 04
% ============================================================


% ------------------------------------------------------------
% SECTION 05
% ------------------------------------------------------------


% ============================================================
% SECTION 05
% INVARIANCE AND CONTINUITY
% Chapter 02: Invariance
% Book I: Ontology of Viability
% ============================================================

\section{Invariance and Continuity}
\label{sec:invariance-and-continuity}

% ------------------------------------------------------------
% 5.1. THE PROBLEM OF CONTINUITY
% ------------------------------------------------------------

The relationship between invariance and continuity
constitutes a further dimension of the ontological
problem developed in this chapter. Both concepts
concern the question of what remains identifiable
when an existent is considered under different
conditions. Nevertheless, they address distinct
questions and must not be treated as interchangeable.

Invariance, according to the provisional relational
definition adopted in this first edition, concerns
the relation between an existent, or one of its
determinate aspects, and a specified domain of
changes or transformations, in relation to which
something remains unchanged in a particular respect.

Continuity, by contrast, concerns whether and how
an existent continues through different conditions,
phases, or states. It raises a question about the
continuation of an existent rather than merely
the unchanged persistence of one of its determinations.

The distinction is important because the continuity
of an existent cannot be established simply by
identifying an invariant feature. Likewise, the
absence of an identified invariant feature does
not automatically establish the discontinuity
of an existent.

The present section examines this distinction
without presupposing a final account of persistence,
temporal identity, or the ontological ground of
continuity.

% ------------------------------------------------------------
% 5.2. CONTINUITY AS A DISTINCT QUESTION
% ------------------------------------------------------------

Continuity is often associated with the continuation
of something across different conditions or phases.
However, this association must be examined carefully.

To say that an existent continues is not necessarily
to say that every determination associated with it
remains unchanged. Different determinations may
undergo alteration while the question of the
continuity of the existent remains meaningful.

The concept of continuity therefore cannot be
reduced to the persistence of an unchanged
configuration or collection of features.

At the same time, the concept cannot be used
without qualification. A claim that an existent
continues requires some account of what is meant
by continuation and what distinguishes it from
the mere succession of different conditions.

The present analysis does not assume that
continuity is necessarily temporal in every
possible ontological context. It initially
examines continuity in relation to different
conditions or phases while leaving the broader
ontological scope of the concept open.

This restriction prevents the analysis from
imposing a temporal model on every form of
ontological continuity before the relevant
questions have been established.

% ------------------------------------------------------------
% 5.3. INVARIANCE DOES NOT ESTABLISH CONTINUITY
% ------------------------------------------------------------

The persistence of a particular determination
under a specified transformation does not,
by itself, establish the continuity of the
existent to which that determination belongs.

An invariant feature may be identified across
different conditions, but this identification
does not automatically explain whether the
existent itself continues through those
conditions.

The distinction can be expressed through
two separate questions:

\begin{enumerate}
    \item Does a particular determination
    remain unchanged under a specified
    domain of transformations?

    \item Does the existent continue through
    the conditions or phases being considered?
\end{enumerate}

The first question concerns invariance.
The second concerns continuity.

Although the answers may be related, neither
answer can be inferred from the other without
an additional argument.

For example, a particular structural relation
may remain unchanged across different
configurations. This may provide relevant
information about the transformations under
consideration, but it does not by itself
establish the ontological continuity of the
existent.

The relation between the invariant
determination and the continuity of the
existent requires further examination.

It would therefore be premature to define
continuity simply as the persistence of an
invariant feature.

% ------------------------------------------------------------
% 5.4. CONTINUITY DOES NOT REQUIRE THE
% INVARIANCE OF EVERY DETERMINATION
% ------------------------------------------------------------

If continuity were understood as the
unchanged persistence of every determination,
it would become difficult to account for
the possibility of continuation through
transformation.

An existent may undergo changes in
configuration, position, quantity, relation,
or other determinations. Such changes do
not automatically settle the question
of whether the existent continues.

The analysis must distinguish between
the continuation of the existent and
the persistence of each particular
determination associated with it.

This distinction does not establish that
an existent necessarily continues through
every possible transformation. It only
shows that the alteration of a particular
determination is not, by itself, sufficient
to establish discontinuity.

The conditions under which continuity
is maintained, altered, or interrupted
remain to be investigated.

Nor should the possibility of continuity
through change be used to introduce an
unexamined assumption about an immutable
substrate or an unchanging essence.

Such explanations may be considered in
comparative philosophical analysis, but
they cannot be adopted as foundational
principles without further justification.

% ------------------------------------------------------------
% 5.5. CONTINUITY AND THE SUCCESSION OF
% CONDITIONS
% ------------------------------------------------------------

A common way of examining continuity is
to consider an existent under successive
conditions.

For analytical purposes, these conditions
may be represented schematically as:

\[
A(t_1), \quad A(t_2), \quad A(t_3)
\]

The notation indicates the consideration
of an existent under different temporal
conditions. It does not, by itself,
establish that the same existent persists
through all of them.

The question of continuity remains distinct
from the formal representation of succession.

A sequence of descriptions may refer to
successive observations, successive states,
or successive conditions. The existence
of such a sequence does not automatically
establish the ontological continuity of
the subject being described.

The analysis must therefore distinguish
between:

\begin{enumerate}
    \item The succession of observations.
    \item The succession of described states.
    \item The continuation of an existent
    through different conditions.
\end{enumerate}

These are related but distinct matters.

The representation of temporal succession
may assist the analysis of continuity,
but it cannot substitute for an ontological
account of what it means for an existent
to continue.

The present chapter does not determine
whether time is the ultimate framework
of continuity or whether continuity
may have forms that cannot be reduced
to temporal succession.

% ------------------------------------------------------------
% 5.6. CONTINUITY, CHANGE, AND INVARIANCE
% ------------------------------------------------------------

The three concepts examined in this
chapter must now be distinguished
within a common analytical framework.

Change concerns the alteration of
one or more determinations under
specified conditions.

Invariance concerns the persistence
of a specified determination in
relation to a defined domain of
changes or transformations.

Continuity concerns whether and
how an existent continues through
different conditions or phases.

These concepts can be related without
being reduced to one another.

A possible case may involve change
in one determination, invariance in
another, and an unresolved question
about the continuity of the existent
as a whole.

Another case may involve continuity
without an immediately identifiable
invariant determination. Whether
such a case is ontologically possible
in a particular domain requires
further investigation.

The analytical distinctions can
be summarized as follows:

\[
\begin{aligned}
\text{Change} &\rightarrow
\text{Alteration of a determination} \\
\text{Invariance} &\rightarrow
\text{Unchanged determination within a domain} \\
\text{Continuity} &\rightarrow
\text{Continuation of an existent}
\end{aligned}
\]

These expressions are conceptual
summaries rather than exhaustive
definitions.

Their purpose is to prevent the
analysis from assuming that the
presence or absence of one of these
dimensions automatically determines
the status of the others.

% ------------------------------------------------------------
% 5.7. CONTINUITY AND IDENTITY
% ------------------------------------------------------------

The relationship between continuity
and identity is distinct from the
relationship between invariance
and identity.

The preceding chapter established
that being and identity are
ontologically inseparable. The
present section does not revise
that commitment.

However, the inseparability of being
and identity does not provide an
immediate explanation of how an
existent continues through different
conditions.

A claim of continuity concerns the
continuation of an existent, while
a claim about identity concerns
the identity-bearing character
of its existence.

These questions are related, but
they cannot be collapsed into a
single proposition.

It would be premature to conclude
that continuity is simply the
temporal persistence of identity,
because such a conclusion would
presuppose an account of the
relationship between identity
and temporal transformation.

It would also be premature to
conclude that continuity is
constituted by the persistence
of invariant features.

The relationship among identity,
continuity, and invariance must
therefore remain open to further
ontological development.

% ------------------------------------------------------------
% 5.8. EPISTEMIC RECOGNITION OF CONTINUITY
% ------------------------------------------------------------

The recognition of continuity
often involves the identification
of relations between different
conditions or phases.

A knower may compare descriptions,
observations, or records and
conclude that an existent has
continued through the conditions
being considered.

Such recognition depends upon
the available evidence, the
criteria of identification,
and the conceptual framework
within which the comparison
is conducted.

However, the epistemic recognition
of continuity must not be confused
with the ontological continuity
of the existent.

An existent may continue without
its continuity being recognized
by a particular knower. Conversely,
a mistaken identification may
lead an observer to attribute
continuity where the relevant
ontological conditions have not
been established.

The limitations of recognition
must therefore be distinguished
from the status of being itself.

Knowledge and non-knowledge remain
intertwined within the situated
condition of the knower. The
recognition of some relations
between successive conditions
may coexist with ignorance of
other relevant determinations.

This does not invalidate the
possibility of knowledge. It
requires that claims about
continuity remain proportionate
to the grounds on which they
are established.

% ------------------------------------------------------------
% 5.9. THE UNRESOLVED GROUND OF CONTINUITY
% ------------------------------------------------------------

The distinction between invariance
and continuity clarifies the
problem but does not resolve
its ontological ground.

Several questions remain open:

\begin{itemize}
    \item What makes the continuation
    of an existent possible?

    \item Is continuity an ontological
    determination, a relation between
    successive conditions, or a
    combination of distinct dimensions?

    \item Can continuity occur without
    the persistence of any particular
    determination identified as invariant?

    \item What is the relationship
    between continuity and the
    identity-bearing existence
    of an existent?

    \item Does continuity have the
    same meaning across different
    ontological realms?

    \item How does continuity relate
    to the dynamics of being?
\end{itemize}

The present chapter does not answer
these questions by presupposing
a permanent substance, an immutable
essence, a fixed structure, or
a universal temporal framework.

Such explanations would introduce
ontological commitments that have
not yet been established.

The question of continuity must
therefore remain open for the
subsequent development of the
system, particularly in relation
to the dynamics of being.

% ------------------------------------------------------------
% 5.10. SYNTHESIS AND TRANSITION
% ------------------------------------------------------------

The analysis has established that
invariance and continuity address
distinct but related ontological
problems.

Invariance concerns what remains
unchanged in a specified respect
within a defined domain of
transformations.

Continuity concerns whether and
how an existent continues through
different conditions or phases.

The persistence of an invariant
determination may be relevant
to the examination of continuity,
but it does not, by itself,
establish continuity.

Likewise, the continuation of
an existent does not necessarily
require the unchanged persistence
of every determination associated
with it.

The distinction prevents two
reductions:

\begin{enumerate}
    \item Continuity must not be
    reduced to the persistence
    of invariant features.

    \item Invariance must not be
    treated as a sufficient
    explanation of the continuation
    of an existent.
\end{enumerate}

The analysis also preserves the
foundational distinction between
identity and continuity. The
identity-bearing character of
existence does not, by itself,
settle the conditions under
which continuity is maintained
through transformation.

The next section examines the
relationship between invariance
and structure. This transition
is necessary because structural
relations may provide important
contexts in which invariance
can be identified, but their
persistence must not be
automatically interpreted as
the ultimate ground of either
identity or continuity.

% ============================================================
% END OF SECTION 05
% ============================================================


% ------------------------------------------------------------
% SECTION 06
% ------------------------------------------------------------


% ============================================================
% SECTION 06
% INVARIANCE AND STRUCTURE
% Chapter 02: Invariance
% Book I: Ontology of Viability
% ============================================================

\section{Invariance and Structure}
\label{sec:invariance-and-structure}

% ------------------------------------------------------------
% 6.1. THE PROBLEM OF STRUCTURAL INVARIANCE
% ------------------------------------------------------------

The relationship between invariance and structure
introduces a further dimension into the analysis
of what remains unchanged under specified
transformations.

An existent may be considered not only in terms
of its individual determinations, but also in
terms of the organization and relations through
which its determinations are connected.

The concept of structure becomes relevant when
the analysis concerns an arrangement of
determinations or relations rather than an
isolated feature.

However, the existence of a structure and the
invariance of that structure are not identical
claims. A structure may undergo transformation,
and a particular structural relation may remain
unchanged while other aspects of the structure
are altered.

The central question of this section is therefore
not whether structure is inherently invariant,
but under what conditions a structure, or a
particular aspect of it, can be described as
invariant.

This inquiry must also avoid assuming that
structure is necessarily the ultimate ground
of identity, continuity, or the existence of
the entity whose organization is being examined.

% ------------------------------------------------------------
% 6.2. STRUCTURE AS AN ORGANIZATION OF
% DETERMINATIONS AND RELATIONS
% ------------------------------------------------------------

For the purposes of the present analysis,
structure refers provisionally to an organized
configuration of determinations and relations
within an existent or within a domain under
consideration.

This formulation is intentionally broad.
It does not presuppose that every structure
must consist of material parts, nor that
every organization must be spatial or
mechanically arranged.

The concept is used here to identify an
organization in which determinations are
considered in relation to one another.

A structure therefore cannot be adequately
described merely by listing the elements
that are present. The relations among
those elements may be constitutive of
the organization under examination.

Nevertheless, the present definition is
analytical rather than a complete ontological
theory of structure.

It does not establish whether structures
are independent ontological realities,
whether they are grounded in the entities
they organize, or whether they emerge
from more fundamental determinations.

Those questions require a broader account
of the realms and dynamics of being.

% ------------------------------------------------------------
% 6.3. STRUCTURE AND STRUCTURAL DETERMINATION
% ------------------------------------------------------------

A distinction must be maintained between
a structure and a particular structural
determination.

A structure refers to an organized
configuration, whereas a structural
determination refers to a particular
feature or relation within that
configuration.

This distinction is important because
the invariance of one structural
determination does not necessarily
entail the invariance of the entire
structure.

For example, a relation between two
elements may remain unchanged while
other relations within the same
organization are modified.

In such a case, it may be justified
to describe the particular relation
as invariant within the specified
domain of transformations.

It would not necessarily be justified
to describe the whole structure
as invariant.

Conversely, a structure may preserve
a specified organization even though
some of its individual elements or
determinations undergo alteration.

The distinction between structural
invariance and the invariance of
particular structural determinations
must therefore remain explicit.

% ------------------------------------------------------------
% 6.4. STRUCTURAL INVARIANCE AND THE
% DOMAIN OF TRANSFORMATION
% ------------------------------------------------------------

As established in the initial definition
of invariance, every claim of invariance
requires a specified domain of
transformations.

This requirement applies equally
to structural invariance.

A structure cannot be described
as invariant without specifying
the transformations in relation
to which its organization remains
unchanged.

The same structure may be invariant
under one domain of transformations
and non-invariant under another.

For example, a particular organization
may remain unchanged under transformations
that alter the position of its elements,
while failing to remain unchanged under
transformations that alter the relations
among those elements.

The difference between these cases
does not constitute a contradiction.
The claims concern different domains
of transformation or different
respects of invariance.

Accordingly, the statement that
a structure is invariant must
always be understood as a qualified
claim.

It must identify, at least analytically:

\begin{enumerate}
    \item The structure or structural
    determination under consideration.

    \item The domain of transformations
    being examined.

    \item The respect in which the
    organization is said to remain
    unchanged.
\end{enumerate}

Without these specifications,
the assertion of structural
invariance remains indeterminate.

% ------------------------------------------------------------
% 6.5. STRUCTURAL INVARIANCE AND CHANGE
% ------------------------------------------------------------

Structural invariance does not imply
the absence of structural change
in every respect.

An organization may preserve certain
relations while undergoing alterations
in other relations, elements, or
configurations.

This possibility requires the
distinction between the invariance
of a specified structural aspect
and the invariance of the structure
considered as a whole.

Let \(S\) represent a structure,
\(T\) a transformation, and
\(R\) a specified structural
determination.

The following expression represents
the invariance of that determination:

\[
R(T(S)) = R(S)
\]

The expression does not establish
that the complete structure remains
unchanged. It only represents the
persistence of the specified
determination \(R\) under \(T\).

A different structural determination,
represented by \(Q\), may undergo
alteration under the same transformation:

\[
Q(T(S)) \neq Q(S)
\]

These expressions illustrate how
structural invariance and structural
change may coexist without contradiction.

The relevant question is always
whether the same determination
is being considered under the
same transformation and in the
same respect.

% ------------------------------------------------------------
% 6.6. STRUCTURE, RELATION, AND ORGANIZATION
% ------------------------------------------------------------

The analysis of structure necessarily
raises the question of the relationship
between structural organization and
the relations that constitute or
characterize it.

A structure is not simply identical
to an arbitrary collection of
relations. The organization in
which relations are situated
matters to the description of
the structure.

At the same time, it would be
premature to assume that every
structure can be reduced entirely
to its relations.

The ontological status of relations
and their relationship to the
entities involved remain broader
questions that cannot be resolved
within the present discussion.

For the analysis of invariance,
the important point is that a
structural determination may
concern a relation rather than
an isolated property.

Consequently, structural invariance
may involve the persistence of
an organization, a relation, or
a specified pattern of relations.

These possibilities must not be
treated as equivalent without
further specification.

In particular, the invariance
of a relation does not necessarily
establish the invariance of the
organization in which that relation
is situated.

% ------------------------------------------------------------
% 6.7. STRUCTURAL INVARIANCE AND IDENTITY
% ------------------------------------------------------------

Structural invariance may be relevant
to the examination of identity,
but it must not be treated as
the ontological foundation of
identity without further argument.

A structure may preserve a particular
organization under a specified
domain of transformations. This
may provide grounds for identifying
a structural continuity or a
particular invariant aspect.

However, the persistence of that
organization does not, by itself,
explain why the existent possesses
the identity it possesses.

The foundational commitment of
this book remains that being
and identity are ontologically
inseparable.

Structure must therefore be examined
in relation to identity-bearing
existence rather than treated as
an independently established
replacement for it.

Nor does the alteration of a
structure automatically establish
the loss of identity.

Such a conclusion would require
an account of the relationship
between the relevant structural
determinations and the identity
of the existent.

The present section does not
presuppose that identity is
constituted by structural
invariance, nor that every
structural transformation
preserves identity.

Both possibilities require
further ontological analysis.

% ------------------------------------------------------------
% 6.8. STRUCTURAL INVARIANCE AND
% EPISTEMIC DESCRIPTION
% ------------------------------------------------------------

Structural invariance is often
identified through comparison.

A knower examines different
configurations and attempts
to determine whether particular
relations or organizational
features remain unchanged.

This process involves the
selection of relevant aspects,
the specification of transformations,
and the construction of descriptions
that make comparison possible.

The resulting claim of structural
invariance is therefore epistemically
mediated.

However, epistemic mediation
does not imply that the structure
or its invariant determinations
are created by the act of recognition.

The distinction between the
ontological status of a structure
and the epistemic identification
of its invariance must be preserved.

A structural relation may remain
unchanged without being recognized
by a particular knower.

Conversely, an apparent structural
invariance may result from an
inadequate description, an
incomplete domain of observation,
or the omission of relevant
determinations.

The knower is situated within
a field of both knowledge and
non-knowledge. The identification
of one structural feature may
coexist with ignorance concerning
other features or relations that
could alter the interpretation.

This limitation does not eliminate
the possibility of knowledge.
It requires that structural
invariance claims remain
appropriately qualified.

% ------------------------------------------------------------
% 6.9. THE LIMITS OF STRUCTURAL INVARIANCE
% ------------------------------------------------------------

The analysis of structural
invariance must remain limited
by the distinctions established
throughout this chapter.

First, the invariance of a
structural determination does
not establish the invariance
of the entire structure.

Second, the invariance of a
structure under one domain
of transformations does not
establish its invariance under
every possible transformation.

Third, structural invariance
does not, by itself, establish
the continuity or identity
of the existent.

Fourth, structural change does
not automatically establish
the loss of identity or the
discontinuity of the existent.

Fifth, the identification of
structural invariance must
not be confused with the
ontological production of
the structure by the knower.

Finally, the present analysis
does not determine whether
structure is ontologically
fundamental, derivative, or
dependent upon more basic
determinations.

These limitations are not
merely restrictions on the
use of terminology. They
identify questions that
require further ontological
development.

The relationship between
structure and the deeper
organization of being must
therefore remain open.

% ------------------------------------------------------------
% 6.10. SYNTHESIS AND TRANSITION
% ------------------------------------------------------------

The present section has examined
the relationship between invariance
and structure.

Structure has been provisionally
understood as an organized
configuration of determinations
and relations.

Structural invariance concerns
the persistence of a specified
structural aspect or organization
under a defined domain of
transformations.

The analysis has established
several distinctions:

\begin{enumerate}
    \item Structure is not identical
    to a particular structural
    determination.

    \item The invariance of one
    structural determination does
    not necessarily establish
    the invariance of the entire
    structure.

    \item Structural invariance
    is always relative to a
    specified domain of
    transformations and a
    particular respect.

    \item Structural invariance
    may coexist with change
    in other structural
    determinations.

    \item Structural invariance
    does not, by itself,
    establish identity,
    continuity, or the
    ontological ground
    of an existent.
\end{enumerate}

The analysis has deliberately
avoided treating structure as
an immutable foundation or
as a complete explanation
of identity-bearing existence.

The next section examines
different forms and domains
of invariance. This examination
will clarify how the relational
definition of invariance can
be applied across different
kinds of determinations without
assuming that all instances
of invariance possess the
same ontological significance.

% ============================================================
% END OF SECTION 06
% ============================================================


% ------------------------------------------------------------
% SECTION 07
% ------------------------------------------------------------


% ============================================================
% SECTION 07
% FORMS AND DOMAINS OF INVARIANCE
% Chapter 02: Invariance
% Book I: Ontology of Viability
% ============================================================

\section{Forms and Domains of Invariance}
\label{sec:forms-and-domains-of-invariance}

% ------------------------------------------------------------
% 7.1. THE NEED TO DISTINGUISH FORMS OF INVARIANCE
% ------------------------------------------------------------

The preceding sections have established a
provisional relational definition of invariance.
Invariance concerns the relation between an
existent, or a determinate aspect of an existent,
and a specified domain of changes or
transformations, within which that existent
or aspect remains unchanged in a particular
respect.

This definition identifies the general structure
of an invariance claim. It does not, however,
imply that every instance of invariance has
the same form or operates within the same
domain.

Different determinations may remain unchanged
under different kinds of transformations.
A relation may remain invariant while a
configuration changes. A quantitative
determination may remain constant while
other determinations vary. A particular
organization may persist under transformations
that alter its elements or their positions.

These differences make it necessary to
distinguish forms and domains of invariance.

The purpose of this section is not to establish
a final and exhaustive classification. Rather,
it is to identify several analytical distinctions
that help specify what is being described
when invariance is attributed to an existent
or one of its determinations.

The classification developed here remains
provisional and does not imply that the
distinguished forms constitute independent
ontological categories.

% ------------------------------------------------------------
% 7.2. THE DOMAIN OF TRANSFORMATION
% ------------------------------------------------------------

The domain of transformation specifies the
changes or transformations in relation to
which invariance is being examined.

Without such a domain, the statement that
something is invariant remains incomplete.

A determination may remain unchanged under
one set of transformations while changing
under another. Therefore, invariance cannot
be attributed without qualification.

For analytical purposes, a domain of
transformation may be specified according
to the kind of alteration being considered,
the conditions under which the alteration
occurs, or the range of transformations
included in the examination.

These specifications are methodological
requirements. They do not imply that
the domain of transformations exists as
an independent ontological entity.

The distinction can be represented
schematically as follows:

\[
\mathcal{D}_1 \neq \mathcal{D}_2
\]

where \(\mathcal{D}_1\) and \(\mathcal{D}_2\)
represent different domains of transformation.

If a determination remains unchanged
within \(\mathcal{D}_1\), this does not
establish that it remains unchanged
within \(\mathcal{D}_2\).

Accordingly, every meaningful claim
of invariance must make its domain
explicit or establish it through
the context of the analysis.

% ------------------------------------------------------------
% 7.3. INVARIANCE OF QUANTITATIVE
% DETERMINATIONS
% ------------------------------------------------------------

One possible form of invariance concerns
the persistence of a quantitative
determination under specified transformations.

A quantity may remain unchanged while
other aspects of the existent undergo
alteration.

For example, a particular quantity
may remain constant under transformations
that modify the configuration, position,
or relations of the elements involved.

This does not establish that the entire
existent remains unchanged. It identifies
only the invariance of the specified
quantitative determination.

The relevant claim can be expressed
schematically as:

\[
Q(T(A)) = Q(A)
\]

where \(Q\) represents a specified
quantitative determination and \(T\)
a transformation within the relevant
domain.

The expression is meaningful only when
the quantity and the transformation
are adequately specified.

It does not establish that quantitative
invariance is more fundamental than
other forms of invariance, nor that
every quantity can remain invariant
under transformation.

The ontological significance of
quantitative determinations remains
dependent upon the broader account
of the existent in which they occur.

% ------------------------------------------------------------
% 7.4. INVARIANCE OF QUALITATIVE
% DETERMINATIONS
% ------------------------------------------------------------

Invariance may also concern qualitative
determinations.

A qualitative determination refers here
to a specified character or feature
of an existent that is not being
considered primarily as a numerical
quantity.

Such a determination may remain
unchanged under some transformations
while other qualitative or quantitative
aspects are modified.

However, the identification of
qualitative invariance requires
careful specification.

A qualitative description may group
together different determinations
under a common term. The persistence
of that description does not
necessarily establish that the
underlying determination has remained
unchanged in every relevant respect.

The analysis must therefore distinguish
between the invariance of a determination
and the continued applicability of
a particular description.

The latter may reflect the limits
or conventions of the descriptive
framework rather than the complete
ontological status of the existent.

This distinction is particularly
important when qualitative descriptions
are used to compare different
conditions or phases.

% ------------------------------------------------------------
% 7.5. RELATIONAL INVARIANCE
% ------------------------------------------------------------

A further form of invariance concerns
relations.

A relation between determinations
or entities may remain unchanged
under transformations that alter
other aspects of the configuration.

For example, the relative ordering
of elements may remain unchanged
even when their positions or
distances are modified.

In such a case, the invariance
concerns the specified relation,
not necessarily the elements,
their properties, or the complete
configuration in which they occur.

Relational invariance must therefore
be distinguished from the invariance
of the relata and from the invariance
of the total relational organization.

The following expression represents
a relation \(R\) that remains
unchanged under a specified
transformation:

\[
R(T(A)) = R(A)
\]

This expression is schematic.
It does not determine the ontological
status of the relation or establish
whether the relation is independent
of the entities involved.

The purpose is only to represent
the possibility that a relation
may be invariant while other
determinations undergo change.

% ------------------------------------------------------------
% 7.6. STRUCTURAL INVARIANCE
% ------------------------------------------------------------

Structural invariance concerns
the persistence of a specified
organization or structural
determination under a defined
domain of transformations.

As discussed in the preceding
section, the invariance of one
structural relation does not
necessarily establish the
invariance of the entire structure.

Likewise, a structure may preserve
a specified organization even
when some of its constituent
determinations are altered.

Structural invariance therefore
requires an explicit distinction
between the organization considered
as a whole and the particular
features or relations through
which that organization is described.

This form of invariance may be
relevant to the examination of
continuity and identity, but
it cannot independently establish
either.

The classification of structural
invariance as a distinct analytical
form does not imply that structure
is an ontologically fundamental
category.

Its deeper status remains open
to the broader development of
the ontology.

% ------------------------------------------------------------
% 7.7. TEMPORAL AND NON-TEMPORAL
% DOMAINS OF INVARIANCE
% ------------------------------------------------------------

The domain of invariance may also
be considered in relation to
the presence or absence of
temporal transformation.

In a temporal domain, a determination
is examined across different
times or successive conditions.

A particular feature may remain
unchanged during a specified
interval, while changing before
or after that interval.

Such invariance is necessarily
qualified by the temporal range
under consideration.

However, not every transformation
must be understood as a temporal
change.

A comparison may concern different
configurations, perspectives,
relations, or conditions without
the analysis being primarily
about temporal succession.

For this reason, the domain
of invariance should not be
reduced in advance to time.

This distinction does not establish
a complete theory of non-temporal
invariance. It only prevents
the assumption that every
invariance claim must be
formulated as persistence
through time.

The ontological relationship
between time, transformation,
and invariance remains open.

% ------------------------------------------------------------
% 7.8. LOCAL AND GENERAL INVARIANCE
% ------------------------------------------------------------

A further distinction concerns
the scope of the determination
being examined.

Local invariance refers to the
persistence of a specified
determination within a limited
aspect, region, relation, or
domain of an existent.

General invariance refers to
a broader claim concerning
a more extensive organization
or range of determinations.

These terms describe differences
in analytical scope rather than
independent ontological categories.

A local invariant does not
necessarily establish a general
invariant.

Likewise, a claim of general
invariance requires more than
the identification of one or
several local invariants.

The distinction can be represented
schematically:

\[
I_{\mathrm{local}}
\not\Rightarrow
I_{\mathrm{general}}
\]

The expression indicates that
local invariance alone does not
logically establish general
invariance.

The same requirement applies
to claims concerning an entire
existent. The invariance of
a limited number of determinations
must not be generalized to
the whole without sufficient
grounds.

% ------------------------------------------------------------
% 7.9. INTERDEPENDENCE AND OVERLAPPING
% DOMAINS
% ------------------------------------------------------------

The forms of invariance distinguished
in this section should not be
understood as mutually exclusive.

A particular determination may
be quantitative and structural,
relational and temporal, or
qualitative and local.

The distinctions identify
different analytical aspects
of invariance rather than
separate classes that must
remain isolated.

Moreover, different domains
of transformation may overlap.

A transformation may alter
several determinations at once,
while preserving others.
Consequently, the identification
of an invariant feature may
require attention to more
than one domain or analytical
dimension.

The classification must therefore
remain flexible enough to
accommodate cases in which
different forms of invariance
are simultaneously relevant.

This flexibility does not
eliminate the need for precision.
On the contrary, the more
domains and determinations
are involved, the more important
it becomes to specify the
respect in which invariance
is being asserted.

% ------------------------------------------------------------
% 7.10. SYNTHESIS AND TRANSITION
% ------------------------------------------------------------

This section has distinguished
several analytical forms and
domains of invariance.

These include:

\begin{enumerate}
    \item Quantitative invariance.
    \item Qualitative invariance.
    \item Relational invariance.
    \item Structural invariance.
    \item Invariance within temporal
    and non-temporal domains.
    \item Local and general invariance.
\end{enumerate}

These distinctions do not
constitute an exhaustive
classification, nor do they
establish separate ontological
kinds of invariance.

They clarify that every claim
of invariance must specify
the determination under
consideration, the domain
of transformation, and the
respect in which something
remains unchanged.

The analysis has also established
that different forms of invariance
may overlap and that the
invariance of a limited aspect
cannot automatically be
generalized to the whole existent.

The next section examines
invariance and knowledge.
It will address the distinction
between the ontological status
of what remains unchanged
and the epistemic conditions
under which such invariance
is identified, described,
and justified.

% ============================================================
% END OF SECTION 07
% ============================================================


% ------------------------------------------------------------
% SECTION 08
% ------------------------------------------------------------


% ============================================================
% SECTION 08
% INVARIANCE AND KNOWLEDGE
% Chapter 02: Invariance
% Book I: Ontology of Viability
% ============================================================

\section{Invariance and Knowledge}
\label{sec:invariance-and-knowledge}

% ------------------------------------------------------------
% 8.1. THE EPISTEMIC PROBLEM OF INVARIANCE
% ------------------------------------------------------------

The analysis of invariance necessarily raises
a question concerning the relationship between
what remains unchanged and the knowledge through
which such invariance is identified.

An invariance claim is not merely an assertion
that something remains unchanged. It also
involves the identification of a determination,
the specification of a domain of transformations,
and the recognition that the determination
persists in a particular respect across that
domain.

These operations involve knowledge.

A knower must distinguish the relevant
determination from other aspects of the
existent, identify or specify the transformations
under consideration, and establish the grounds
for claiming that the determination remains
unchanged.

However, the epistemic conditions under which
invariance is recognized must not be confused
with the ontological status of invariance itself.

The fact that invariance is known through
particular acts of observation, comparison,
or reasoning does not establish that it is
produced by those acts.

The central problem of this section is
therefore to distinguish the ontological
question of invariance from the epistemic
conditions through which invariance becomes
knowable.

% ------------------------------------------------------------
% 8.2. ONTOLOGICAL INVARIANCE AND
% EPISTEMIC RECOGNITION
% ------------------------------------------------------------

The distinction between ontological invariance
and epistemic recognition is fundamental.

Ontological invariance concerns the persistence
of a specified determination in relation to
a domain of transformations.

Epistemic recognition concerns the process
through which a knower identifies, describes,
and justifies the claim that such persistence
occurs.

These two dimensions are related, but they
are not identical.

An existent or one of its determinations
may remain unchanged even when no knower
recognizes that invariance.

Conversely, a knower may identify an apparent
invariance that is not adequately supported
by the relevant conditions or evidence.

The possibility of error therefore belongs
to the epistemic examination of invariance.
It does not imply that the ontological status
of what is being examined is reducible to
the knower's recognition.

The distinction can be summarized as follows:

\begin{enumerate}
    \item Ontological invariance concerns
    what remains unchanged under specified
    transformations.

    \item Epistemic recognition concerns
    how the persistence of that determination
    is identified and justified.
\end{enumerate}

Neither dimension should be substituted
for the other.

% ------------------------------------------------------------
% 8.3. THE CONDITIONS OF RECOGNITION
% ------------------------------------------------------------

The recognition of invariance requires
several analytical operations.

First, the knower must identify the
existent or determinate aspect under
consideration.

Second, the relevant domain of
transformations must be specified.
Without such specification, the
claim of invariance lacks a clear
scope.

Third, the respect in which the
determination is said to remain
unchanged must be identified.

Fourth, the available grounds for
the comparison must be examined.

These operations may involve
observation, measurement, conceptual
analysis, formal reasoning, or
other epistemic procedures.

Their relative importance depends
upon the nature of the existent
and the determination being examined.

The conditions of recognition
are therefore not necessarily
identical across different forms
and domains of invariance.

A quantitative invariance claim
may require different grounds
from a claim concerning a
qualitative, relational, or
structural determination.

This does not imply that one
form of knowledge is universally
superior to another. It indicates
that the adequacy of an epistemic
procedure must be assessed in
relation to the specific claim
being made.

% ------------------------------------------------------------
% 8.4. COMPARISON AND THE IDENTIFICATION
% OF INVARIANCE
% ------------------------------------------------------------

The identification of invariance
often involves comparison between
different conditions or configurations.

A knower examines what is present
under one condition and compares
it with what is present under
another.

The comparison seeks to establish
whether a specified determination
has remained unchanged.

However, comparison does not
automatically establish invariance.

The conditions being compared
may differ in ways that are not
recognized. The descriptions
may omit relevant determinations.
The transformations may not be
adequately specified.

Furthermore, two descriptions
may appear identical while
referring to different ontological
conditions.

The epistemic act of comparison
must therefore be distinguished
from the ontological relation
of invariance that it seeks
to identify.

A comparison is an operation
of the knower. Invariance is
the relation being examined.

The adequacy of the comparison
depends upon whether it captures
the relevant determination and
the transformations within
the specified domain.

% ------------------------------------------------------------
% 8.5. MEASUREMENT, DESCRIPTION,
% AND THE LIMITS OF REPRESENTATION
% ------------------------------------------------------------

In some cases, the identification
of invariance involves measurement
or formal description.

Such procedures can make it
possible to compare determinations
across different conditions
with greater precision.

Nevertheless, the result of
a measurement must not be
identified without qualification
with the complete ontological
status of the existent.

Every measurement operates
through particular procedures,
units, instruments, distinctions,
and assumptions.

A measured quantity may remain
constant within the resolution
and conditions of a procedure,
while other changes remain
unidentified.

Likewise, a formal description
may preserve a particular
relation while omitting other
aspects of the existent.

The limits of measurement and
representation must therefore
be included in the assessment
of an invariance claim.

This does not make measurement
or formalization epistemically
invalid. It establishes that
their results must be interpreted
within the scope of the procedures
through which they are obtained.

The distinction between a
determination and its representation
must remain explicit.

% ------------------------------------------------------------
% 8.6. KNOWLEDGE AND NON-KNOWLEDGE
% ------------------------------------------------------------

The epistemic examination of
invariance takes place within
the situated condition of a
knower.

The knower does not occupy
a position from which every
determination of an existent
is simultaneously available.

Knowledge of one aspect may
coexist with non-knowledge
of other aspects. The identification
of one invariant determination
may occur while other relevant
determinations remain unknown.

These conditions are not
necessarily independent or
sequential.

The knowledge of a particular
domain and the non-knowledge
of other domains may be
intertwined within the same
epistemic situation.

Consequently, the recognition
of invariance must not be
interpreted as the achievement
of complete knowledge of the
existent.

Nor should the limits of
knowledge be treated as evidence
that invariance does not exist.

The absence of an identified
invariant is not equivalent
to the demonstrated absence
of invariance.

Likewise, the recognition of
an invariant in one domain
does not establish that the
existent is invariant in every
other domain.

The epistemic position of the
knower must therefore remain
explicit in every sufficiently
qualified claim about invariance.

% ------------------------------------------------------------
% 8.7. THE DANGER OF EPISTEMIC
% REDUCTION
% ------------------------------------------------------------

A reduction of invariance to
epistemic recognition would
make the existence of invariance
dependent upon whether a knower
has identified it.

Such a position would confuse
the conditions of knowing
with the status of what is
known.

The ontological analysis adopted
in this book does not accept
this reduction.

The existent is not constituted
by the knower's description
of it, and the identity-bearing
character of being is not
exhausted by the knowledge
available to a particular
subject.

Accordingly, invariance must
be examined as a relation
concerning an existent or
one of its determinations,
not merely as a feature of
an epistemic representation.

At the same time, rejecting
epistemic reduction does not
justify ignoring the conditions
and limitations of knowledge.

An ontological claim that
is not epistemically justified
cannot be presented as an
established result merely
because it is asserted to
concern reality.

The analysis must therefore
preserve both dimensions:
the ontological independence
of the existent from a
particular act of knowing,
and the epistemic responsibility
required to make justified
claims about it.

% ------------------------------------------------------------
% 8.8. INVARIANCE, KNOWLEDGE,
% AND THE SPECIFICATION OF DOMAINS
% ------------------------------------------------------------

The specification of a domain
of transformations is both
an analytical requirement
of invariance and an epistemic
requirement for recognizing it.

The domain determines the
range within which the
invariance claim is being
examined.

The knower's ability to identify
that domain, however, may be
limited by the available
knowledge and the conceptual
resources through which
transformations are described.

A domain may be incompletely
specified, or a transformation
may be overlooked.

In such cases, an invariance
claim may be valid within
the transformations actually
examined but unjustifiably
extended beyond them.

The distinction between
a specified domain and an
assumed universal domain
is therefore essential.

The recognition of invariance
within a limited domain does
not establish invariance
under every possible transformation.

This limitation applies even
when the identified invariance
is strongly supported within
the domain that has been
examined.

The epistemic justification
of an invariance claim must
therefore remain proportionate
to the scope of the claim.

% ------------------------------------------------------------
% 8.9. THE STATUS OF INVARIANCE
% AS AN OPEN ONTOLOGICAL QUESTION
% ------------------------------------------------------------

The relational definition of
invariance adopted in this
first edition provides a
framework for distinguishing
what remains unchanged under
specified transformations.

It also makes possible a
disciplined examination of
the relationship between
invariance, change, identity,
continuity, and structure.

However, the adequacy of this
definition as a complete
ontological account has not
been established.

It remains possible that
invariance has a deeper
ontological significance
that cannot be fully expressed
through the relational
definition currently employed.

This possibility must not
be treated as an established
conclusion.

The question is intentionally
left open for later examination,
when the broader ontological
framework has been developed
sufficiently to reconsider
the status of invariance.

The distinction between
ontological invariance and
epistemic recognition remains
necessary regardless of whether
the relational definition is
ultimately retained or revised.

A deeper ontological account
would still need to distinguish
what invariance is from the
conditions under which it
becomes known.

% ------------------------------------------------------------
% 8.10. SYNTHESIS AND TRANSITION
% ------------------------------------------------------------

This section has examined
the relationship between
invariance and knowledge.

The analysis has established
that the ontological status
of invariance must be
distinguished from the
epistemic processes through
which invariance is recognized.

The identification of invariance
involves the specification
of a determination, a domain
of transformations, and the
respect in which something
is said to remain unchanged.

These operations are epistemically
mediated, but they do not
thereby produce the ontological
relation being examined.

The analysis has also emphasized
that knowledge and non-knowledge
are intertwined within the
situated condition of the
knower.

Consequently:

\begin{enumerate}
    \item The absence of knowledge
    of an invariant does not
    establish its ontological
    absence.

    \item The recognition of
    an invariant within one
    domain does not establish
    universal invariance.

    \item The results of
    measurement and description
    must be interpreted within
    the limits of their
    procedures.

    \item The ontological
    independence of invariance
    from a particular act
    of knowing does not
    eliminate the need for
    epistemic justification.

    \item The possibility of
    a deeper ontological
    significance of invariance
    remains open.
\end{enumerate}

The next section examines
the scope and limits of
the concept of invariance.
It will consolidate the
commitments established
throughout this chapter,
distinguish them from
unresolved questions, and
clarify the extent to which
the provisional relational
definition can be used
within the ontology of
viability.

% ============================================================
% END OF SECTION 08
% ============================================================


% ------------------------------------------------------------
% SECTION 09
% ------------------------------------------------------------

% ============================================================
% SECTION 09
% SCOPE AND LIMITS OF INVARIANCE
% Chapter 02: Invariance
% Book I: Ontology of Viability
% ============================================================

\section{Scope and Limits of Invariance}
\label{sec:scope-and-limits-of-invariance}

% ------------------------------------------------------------
% 9.1. THE NEED TO ESTABLISH THE SCOPE OF INVARIANCE
% ------------------------------------------------------------

The preceding sections have developed a provisional
relational account of invariance and examined its
relationship with identity, difference, change,
continuity, structure, different forms of determination,
and knowledge.

These examinations have clarified the analytical
use of invariance, but they have also established
the need to delimit its scope.

A concept cannot be used adequately if the conditions
under which it applies remain unspecified. Invariance
is particularly vulnerable to this problem because
claims of invariance may be extended beyond the
determinations, transformations, and domains for
which they have been established.

The statement that something is invariant is never
sufficient by itself. It must be possible to identify
what is invariant, under which transformations,
and in what respect it remains unchanged.

The present section consolidates these requirements
and distinguishes the commitments adopted in this
chapter from the questions that remain unresolved.

Its purpose is not to establish an exhaustive
ontology of invariance. Rather, it clarifies the
legitimate scope of the concept within the
provisional framework of the first edition.

% ------------------------------------------------------------
% 9.2. THE PROVISIONAL RELATIONAL DEFINITION
% ------------------------------------------------------------

For the purposes of this first edition, invariance
is defined as a relation between an existent, or
a determinate aspect of an existent, and a specified
domain of changes or transformations, such that
the existent or aspect remains unchanged in a
particular respect under those changes or
transformations.

This definition establishes three analytical
components:

\begin{enumerate}
    \item \textbf{The subject of invariance:}
    the existent or determinate aspect whose
    invariance is being examined.

    \item \textbf{The domain of transformation:}
    the specified changes or transformations
    in relation to which invariance is asserted.

    \item \textbf{The respect of invariance:}
    the particular determination, feature,
    or relation that remains unchanged.
\end{enumerate}

These components are analytical distinctions
within a single invariance claim. They must
not be interpreted as three independently
existing ontological entities.

The definition is intentionally relational.
It provides a framework for identifying and
examining invariance without presupposing
that invariance is itself an independent
substance, an immutable foundation, or
a universal property of every existent.

The relational definition is a methodological
and conceptual commitment of the present
edition. It does not constitute a final
metaphysical conclusion concerning the
ultimate ontological status of invariance.

% ------------------------------------------------------------
% 9.3. THE REQUIREMENT OF SPECIFICATION
% ------------------------------------------------------------

Every claim of invariance must be sufficiently
specified to distinguish it from an unrestricted
assertion of permanence or immutability.

The subject, domain, and respect of invariance
must be identifiable, either explicitly or
through the context in which the claim is made.

An assertion that a structure remains invariant,
for example, is incomplete unless it is possible
to determine which structural aspect is being
considered and which transformations are
included in the relevant domain.

The same requirement applies to quantitative,
qualitative, relational, and other forms of
invariance.

The scope of an invariance claim must also
remain proportionate to the grounds on which
it is established.

The recognition of invariance within a limited
domain does not justify its extension to
transformations that have not been considered.

Accordingly, the following inference is not
generally justified:

\[
\text{Invariance within } \mathcal{D}_1
\not\Rightarrow
\text{Invariance within every domain}
\]

The expression indicates that the identification
of invariance within one specified domain does
not establish invariance under all possible
transformations.

This limitation is not a weakness of the concept.
It is a condition of its precise application.

% ------------------------------------------------------------
% 9.4. INVARIANCE IS NOT IDENTITY
% ------------------------------------------------------------

The distinction between invariance and identity
is one of the central commitments of this chapter.

Identity concerns the identity-bearing character
of an existent, which, according to the foundational
position of this book, is ontologically inseparable
from being.

Invariance, as provisionally defined here,
concerns the persistence of a specified
determination under a specified domain of
transformations.

The two concepts therefore address different
questions.

The identity of an existent cannot be reduced
to the invariance of one or more of its
determinations.

Likewise, the identification of an invariant
determination does not, by itself, establish
the identity of the existent.

The distinction also prevents a premature
conclusion that every transformation of an
existent must either preserve or destroy
its identity according to a single invariant
criterion.

The relationship between identity and
transformation requires a broader ontological
account than the present chapter provides.

The origin and deeper constitution of identity
remain deferred to the subsequent development
of the realms of being, light, and ultimately
the dynamics of being.

% ------------------------------------------------------------
% 9.5. INVARIANCE IS NOT IMMUTABILITY
% ------------------------------------------------------------

Invariance must also be distinguished from
immutability.

Immutability, in its unrestricted sense,
suggests the absence of change. The relational
concept of invariance developed here does
not require the absence of every possible
change.

An existent may undergo transformations
while a particular determination remains
unchanged within the specified domain.

Consequently, the identification of an
invariant aspect does not establish that
the existent is immutable as a whole.

Nor does the absence of invariance in
one respect establish that the existent
is subject to unrestricted change in
every other respect.

The distinction prevents the concept
of invariance from being used to introduce
an unexamined metaphysics of absolute
permanence.

It also prevents change from being treated
as the universal negation of invariance.

The relevant question remains whether
a particular determination is unchanged
under the transformations being considered.

% ------------------------------------------------------------
% 9.6. INVARIANCE IS NOT CONTINUITY
% ------------------------------------------------------------

The distinction between invariance and
continuity must remain equally explicit.

Invariance concerns the persistence of
a specified determination under a defined
domain of transformations.

Continuity concerns whether and how an
existent continues through different
conditions or phases.

The invariance of a particular feature
may be relevant to an examination of
continuity, but it does not independently
establish that continuity.

Conversely, the continuation of an
existent does not necessarily require
the unchanged persistence of every
determination associated with it.

The concept of continuity therefore
cannot be reduced to the persistence
of invariant features.

This distinction also prevents the
present account from presupposing
a permanent substrate or fixed essence
as the necessary ground of continuation.

The ontological relationship between
identity, continuity, and transformation
remains open for examination within
the broader development of the system.

% ------------------------------------------------------------
% 9.7. INVARIANCE AND STRUCTURE:
% THE LIMIT OF GENERALIZATION
% ------------------------------------------------------------

The analysis of structural invariance
has established that the persistence
of a particular structural determination
does not necessarily entail the invariance
of the entire structure.

A relation may remain unchanged while
other relations or determinations are
altered.

Likewise, a structure may preserve
a specified organization under some
transformations without remaining
unchanged in every respect.

The distinction between local and
general invariance is therefore essential.

The invariance of a limited aspect
must not be generalized to the whole
existent or its entire organization
without adequate grounds.

This limitation applies not only to
structural invariance but also to
quantitative, qualitative, and relational
invariance.

No particular form of invariance
automatically establishes a universal
claim concerning the existent as a whole.

The scope of every claim must remain
tied to the determination and domain
that have actually been examined.

% ------------------------------------------------------------
% 9.8. EPISTEMIC LIMITS AND ONTOLOGICAL STATUS
% ------------------------------------------------------------

The identification of invariance is
epistemically mediated.

A knower must identify the relevant
determination, specify the domain of
transformations, and establish the
grounds for claiming that something
remains unchanged.

These operations may involve observation,
comparison, measurement, formalization,
or conceptual analysis.

However, the epistemic conditions
through which invariance is recognized
must not be confused with its ontological
status.

The absence of knowledge concerning
a particular invariant does not establish
that no such invariance exists.

Similarly, the recognition of invariance
within one domain does not establish
its presence within every other domain.

Knowledge and non-knowledge are
intertwined within the situated condition
of the knower. Knowledge of one aspect
of an existent may coexist with
non-knowledge of other aspects that
could be relevant to the interpretation
of its invariance.

This situation does not justify either
an unrestricted skepticism or an
unqualified claim of knowledge.

It requires that ontological assertions
and epistemic justifications be
distinguished and appropriately
related.

The existent is not reduced to what
a particular knower recognizes, but
claims concerning the existent remain
subject to the epistemic grounds
through which they are justified.

% ------------------------------------------------------------
% 9.9. THE OPEN ONTOLOGICAL STATUS OF INVARIANCE
% ------------------------------------------------------------

The relational definition adopted in
this first edition provides a workable
framework for examining invariance
across different determinations and
domains of transformation.

Nevertheless, the question remains
whether invariance possesses a deeper
ontological significance that cannot
be fully expressed through its
relational formulation.

The present chapter does not resolve
this question.

It would be premature to conclude
that invariance is merely an epistemic
construct, because the relational
definition concerns an existent or
one of its determinations and does
not reduce its status to the act
of knowing.

It would be equally premature to
conclude that invariance constitutes
an independent ontological principle
or a fundamental condition of being.

Neither conclusion has been
established by the analysis developed
so far.

The decision to employ a relational
definition in the first edition
is therefore a provisional
methodological commitment.

Its adequacy may be reconsidered
in subsequent revisions, when
the broader ontological framework
has been developed sufficiently
to examine whether invariance
has a deeper status.

This openness is deliberate.
It preserves the distinction between
what the present framework can
legitimately establish and what
remains a question for further
ontological investigation.

% ------------------------------------------------------------
% 9.10. ESTABLISHED COMMITMENTS AND
% UNRESOLVED QUESTIONS
% ------------------------------------------------------------

The analysis of this chapter
has established a number of
conceptual and methodological
commitments.

These commitments can be
distinguished from questions
that remain unresolved.

\subsection*{Established Commitments}

The following principles govern
the use of invariance within
the present chapter:

\begin{enumerate}
    \item Invariance is provisionally
    defined as a relation between
    an existent or a determinate
    aspect of an existent and
    a specified domain of
    transformations.

    \item Every invariance claim
    requires the specification
    of its subject, domain,
    and respect.

    \item Invariance is not
    identical to identity.

    \item Invariance is not
    equivalent to immutability
    or the absence of all change.

    \item Invariance is not
    identical to continuity.

    \item The invariance of
    a particular determination
    does not automatically
    establish the invariance
    of the whole existent.

    \item Different forms and
    domains of invariance may
    overlap and are not
    necessarily mutually
    exclusive.

    \item The epistemic recognition
    of invariance must be
    distinguished from its
    ontological status.

    \item Knowledge and
    non-knowledge must both
    be considered when the
    epistemic justification
    of invariance claims
    is examined.

    \item The relational
    definition is a provisional
    commitment of the first
    edition rather than a
    final metaphysical
    conclusion.
\end{enumerate}

\subsection*{Unresolved Questions}

The following questions remain
open and must not be treated
as settled conclusions:

\begin{enumerate}
    \item Does invariance possess
    an ontological significance
    beyond its relational
    formulation?

    \item What is the deeper
    relationship between
    invariance and the
    identity-bearing character
    of existence?

    \item Can invariance be
    understood as a condition
    of certain forms of
    continuity, or is its
    relevance always dependent
    upon the specific domain
    being examined?

    \item How does invariance
    relate to the different
    realms of being?

    \item Does the ontological
    status of invariance
    differ across different
    realms or forms of
    manifestation?

    \item What is the relationship
    between invariance and
    the dynamics of being?

    \item Is there a deeper
    ground that explains
    why certain determinations
    remain invariant through
    particular transformations?
\end{enumerate}

These questions identify the
limits of the present account.
They are not reasons to abandon
the relational definition,
but reasons to preserve its
provisional status.

% ------------------------------------------------------------
% 9.10.1. TABLE: SCOPE AND LIMITS OF INVARIANCE
% ------------------------------------------------------------

% ============================================================
% TABLE 09 — SCOPE AND LIMITS OF INVARIANCE
% Book I: Ontology of Viability
% Part I: Identity and Manifestation
% Chapter 2: Invariance
% ============================================================

\begin{table}[htbp]
    \centering
    \caption{The Scope and Limits of the Concept of Invariance}
    \label{tab:scope-and-limits-of-invariance}

    \renewcommand{\arraystretch}{1.35}

    \begin{tabular}{|p{0.23\textwidth}|p{0.34\textwidth}|p{0.34\textwidth}|}
        \hline
        \textbf{Dimension} &
        \textbf{What the Concept Establishes} &
        \textbf{What It Does Not Establish} \\
        \hline

        Subject of invariance &
        An existent, or a determinate aspect of an existent, may be considered with respect to invariance. &
        The subject of invariance is not thereby reduced to the particular determination under consideration. \\
        \hline

        Domain of change &
        Invariance is assessed in relation to a specified domain of changes or transformations. &
        Invariance within one domain does not establish invariance under every possible change or transformation. \\
        \hline

        Respect of invariance &
        A particular determination, feature, or relation may remain unchanged under the specified transformations. &
        The invariance of one respect does not establish the invariance of every aspect of the existent. \\
        \hline

        Identity &
        Invariance may help describe how a determinate aspect remains unchanged under specified transformations. &
        Invariance does not constitute identity, provide its complete explanation, or establish that identity is exhausted by invariant features. \\
        \hline

        Difference &
        Different determinations may coexist with the invariance of a particular determination. &
        Difference does not necessarily negate invariance, and invariance does not eliminate difference. \\
        \hline

        Change &
        An existent may undergo change in some respects while remaining invariant in another respect. &
        Invariance does not mean the total absence of change or the immutability of the existent as a whole. \\
        \hline

        Continuity &
        Invariance may contribute to the description of certain patterns of continuity. &
        The invariance of a particular feature does not, by itself, establish the continuity of the existent as a whole. \\
        \hline

        Structure &
        Particular structural determinations or relations may remain invariant under specified transformations. &
        The invariance of one structural relation does not establish the invariance of the entire structure. \\
        \hline

        Knowledge &
        Invariance may be investigated, recognized, compared, and formally described through epistemic operations. &
        Epistemic recognition does not constitute ontological invariance, and the absence of recognition does not establish its absence. \\
        \hline

        Ontological significance &
        The relational concept provides an initial framework for describing determinate invariance. &
        This framework does not settle whether invariance has a deeper ontological status beyond the relational account adopted in this edition. \\
        \hline

    \end{tabular}
\end{table}

% ============================================================
% END OF TABLE 09
% ============================================================

% ------------------------------------------------------------
% 9.11. SYNTHESIS AND TRANSITION
% ------------------------------------------------------------

The present section has clarified
the scope within which invariance
can be used in the ontology of
viability.

The relational definition provides
a precise analytical framework
for examining what remains
unchanged under specified
transformations.

Its use requires the identification
of the existent or determination
under consideration, the domain
of transformation, and the
particular respect in which
invariance is asserted.

The analysis has also established
that invariance cannot be reduced
to identity, immutability,
continuity, structure, or
epistemic recognition.

These concepts are related,
but each addresses a distinct
problem.

The limits identified throughout
the chapter are essential to
preventing the concept of
invariance from carrying
metaphysical commitments
that have not yet been justified.

The relational definition is
therefore retained as the
working framework of the
first edition, while its
possible deeper ontological
significance remains open
for subsequent reconsideration.

The final section of this chapter
will synthesize the conceptual
results, consolidate the
established distinctions,
and identify the role of
invariance within the broader
development of the ontology
of viability.

% ============================================================
% END OF SECTION 09
% ============================================================

% ------------------------------------------------------------
% SECTION 10
% ------------------------------------------------------------

% ============================================================
% SECTION 10
% SYNTHESIS OF INVARIANCE
% Chapter 02: Invariance
% Book I: Ontology of Viability
% ============================================================

\section{Synthesis of Invariance}
\label{sec:synthesis-of-invariance}

% ------------------------------------------------------------
% 10.1. THE ROLE OF INVARIANCE IN THE ONTOLOGY
% ------------------------------------------------------------

The preceding sections have examined invariance
as a distinct conceptual problem within the
ontology of viability.

The inquiry began with the question of what
remains unchanged when an existent, or one
of its determinations, is considered in relation
to changes and transformations.

The analysis gradually established that
invariance cannot be understood adequately
through an unrestricted assertion of permanence.
It requires the specification of what is being
examined, the transformations under consideration,
and the respect in which something remains
unchanged.

For the purposes of this first edition, invariance
has therefore been provisionally defined as a
relation between an existent, or a determinate
aspect of an existent, and a specified domain
of changes or transformations, such that the
existent or aspect remains unchanged in a
particular respect under those transformations.

This definition provides the working framework
of the present chapter.

It enables the examination of invariance without
prematurely identifying it with identity,
immutability, continuity, structure, or any
other ontological determination.

The purpose of this final section is to consolidate
the distinctions established throughout the
chapter and clarify the conceptual position
that invariance occupies within the broader
development of the ontology of viability.

% ------------------------------------------------------------
% 10.2. THE THREE COMPONENTS OF AN
% INVARIANCE CLAIM
% ------------------------------------------------------------

Every sufficiently specified invariance claim
contains three analytical components.

\begin{enumerate}
    \item \textbf{Subject of invariance:}
    the existent or determinate aspect
    whose invariance is being examined.

    \item \textbf{Domain of transformation:}
    the specified changes or transformations
    in relation to which invariance is asserted.

    \item \textbf{Respect of invariance:}
    the particular determination, feature,
    or relation that remains unchanged.
\end{enumerate}

These components are analytically distinguishable,
but they are not independent ontological entities.
They belong to the specification of a single
invariance claim.

The general structure may be represented
schematically as:

\[
\mathcal{I}(A,\mathcal{D},R)
\]

where:

\begin{itemize}
    \item \(A\) represents the existent or
    determinate aspect under consideration;

    \item \(\mathcal{D}\) represents the
    specified domain of transformations;

    \item \(R\) represents the determination
    or respect in which invariance is asserted.
\end{itemize}

This notation is a conceptual representation
rather than a complete formal definition.
It does not establish a mathematical theory
of invariance or settle its ultimate ontological
status.

Its purpose is to make explicit the conditions
that must be distinguished whenever invariance
is attributed to an existent.

% ------------------------------------------------------------
% 10.3. INVARIANCE AND IDENTITY
% ------------------------------------------------------------

The chapter has consistently distinguished
invariance from identity.

The foundational position of this book is
that being and identity are ontologically
inseparable. Existence is not first established
as an identity-neutral condition to which
identity is subsequently added.

Invariance, however, addresses a different
question.

It concerns whether a specified determination
remains unchanged under a defined domain
of transformations.

The identity of an existent cannot therefore
be reduced to the invariance of one or more
of its determinations.

Likewise, the identification of an invariant
determination does not, by itself, explain
the identity of the existent.

The distinction is essential because the
persistence of a particular feature and
the identity-bearing character of an existent
are not equivalent ontological claims.

The relationship between identity and
transformation remains open for further
development, particularly in relation to
the deeper constitution and dynamics
of being.

% ------------------------------------------------------------
% 10.4. INVARIANCE AND DIFFERENCE
% ------------------------------------------------------------

The analysis of difference established
that difference and invariance are not
necessarily opposed.

Different determinations may coexist
within the same existent, and a transformation
may alter one determination while another
remains unchanged.

Consequently, the existence of difference
does not automatically negate invariance.

At the same time, the same determination
cannot be both changed and unchanged
under the same transformation, in the
same respect, and within the same domain.

This qualification preserves the coherence
of invariance claims.

The distinction between difference and
invariance also prevents the analysis
from assuming that every difference
must produce a loss of identity or that
every invariant determination eliminates
the possibility of difference.

Invariance and difference therefore
identify distinct aspects of the relation
between an existent and its determinations.

% ------------------------------------------------------------
% 10.5. INVARIANCE AND CHANGE
% ------------------------------------------------------------

The relationship between invariance
and change has been established through
the distinction between a specified
unchanged determination and the
alteration of determinations under
transformation.

Invariance is not the absence of all change.

An existent may undergo changes in
some respects while a particular
determination remains invariant
within the relevant domain.

The scope of the invariance claim
must therefore remain limited to
the determination and transformations
actually specified.

Conversely, the identification of
change in one determination does
not establish that every other
determination has changed.

The relationship between invariance
and change is consequently domain-
and determination-dependent.

This formulation prevents both
the reduction of invariance to
absolute permanence and the
reduction of change to the universal
negation of invariance.

% ------------------------------------------------------------
% 10.6. INVARIANCE AND CONTINUITY
% ------------------------------------------------------------

The distinction between invariance
and continuity has shown that
the persistence of a particular
determination does not, by itself,
establish the continuity of an
existent.

Continuity concerns whether and
how an existent continues through
different conditions or phases.

Invariance concerns whether a
specified determination remains
unchanged under a defined domain
of transformations.

An invariant feature may be relevant
to the examination of continuity,
but it cannot automatically serve
as its sufficient explanation.

Likewise, continuity does not
necessarily require the unchanged
persistence of every determination
associated with an existent.

The ontological ground of continuity,
and its relationship to identity
through transformation, remain
open questions.

The present account therefore
preserves the distinction without
assuming that invariance is either
a necessary or a sufficient condition
of continuity.

% ------------------------------------------------------------
% 10.7. INVARIANCE AND STRUCTURE
% ------------------------------------------------------------

The examination of structure
has established that structural
invariance may concern an
organization, a relation, or
a particular structural
determination.

The invariance of one structural
feature does not necessarily
establish the invariance of
the entire structure.

Likewise, the invariance of
an organization under a specified
domain of transformations does
not establish that the structure
is immutable under every possible
transformation.

Structure provides an important
context in which invariance
can be examined, but it must
not be treated as an automatically
established ontological foundation
for identity or continuity.

The deeper ontological status
of structure remains open to
the broader development of
the system.

% ------------------------------------------------------------
% 10.8. FORMS AND DOMAINS OF INVARIANCE
% ------------------------------------------------------------

The chapter has distinguished
several analytical forms of
invariance, including:

\begin{enumerate}
    \item Quantitative invariance.
    \item Qualitative invariance.
    \item Relational invariance.
    \item Structural invariance.
    \item Invariance within temporal
    and non-temporal domains.
    \item Local and general invariance.
\end{enumerate}

These distinctions do not
constitute an exhaustive
classification of invariance.

Nor do they establish mutually
exclusive ontological categories.

Different forms may overlap
within a single case. A particular
determination may be quantitative
and structural, relational and
temporal, or qualitative and local.

The classification is intended
to clarify the different analytical
respects in which invariance
may be examined.

It also reinforces the requirement
that every claim of invariance
must specify its relevant domain
and the determination under
consideration.

% ------------------------------------------------------------
% 10.9. INVARIANCE AND KNOWLEDGE
% ------------------------------------------------------------

The chapter has distinguished
the ontological status of
invariance from the epistemic
conditions through which
it is recognized.

The identification of invariance
involves observation, comparison,
measurement, description,
or reasoning.

These procedures provide
grounds for epistemic claims,
but they do not thereby produce
the ontological relation being
examined.

An existent or one of its
determinations may remain
unchanged without being
recognized by a particular
knower.

Conversely, an apparent
invariance may result from
an incomplete description,
an inadequately specified
domain, or the omission
of relevant determinations.

Knowledge and non-knowledge
are intertwined within the
situated condition of the
knower. The recognition
of one invariant determination
may coexist with non-knowledge
of other aspects that could
be relevant to the analysis.

This does not justify reducing
invariance to knowledge, nor
does it justify ignoring the
epistemic grounds of invariance
claims.

Ontological independence and
epistemic justification must
remain distinct but related
requirements.

% ------------------------------------------------------------
% 10.10. ESTABLISHED COMMITMENTS
% ------------------------------------------------------------

The analysis of this chapter
establishes the following
commitments for the first
edition:

\begin{enumerate}
    \item Invariance is provisionally
    understood as a relation
    between an existent or
    a determinate aspect of
    an existent and a specified
    domain of transformations.

    \item Every invariance claim
    must identify its subject,
    domain of transformation,
    and respect of invariance.

    \item The analytical components
    of an invariance claim
    must not be treated as
    independent ontological
    entities.

    \item Invariance is distinct
    from identity, difference,
    change, continuity, and
    structure, even where
    these concepts are related.

    \item Invariance does not
    imply immutability or
    the absence of all change.

    \item The invariance of
    a particular determination
    does not automatically
    establish the invariance
    of an existent as a whole.

    \item Different forms and
    domains of invariance
    may overlap.

    \item Epistemic recognition
    must be distinguished
    from ontological invariance.

    \item The recognition of
    invariance within a
    limited domain does not
    establish universal
    invariance.

    \item The relational
    definition remains
    provisional and does
    not settle the ultimate
    ontological status of
    invariance.
\end{enumerate}

These commitments delimit
the legitimate use of the
concept within the present
stage of the ontology.

% ------------------------------------------------------------
% 10.11. UNRESOLVED QUESTIONS
% ------------------------------------------------------------

Despite the distinctions
established in this chapter,
several questions remain
unresolved.

First, it remains open
whether invariance has
a deeper ontological
significance beyond its
relational formulation.

Second, the relationship
between invariance and
the identity-bearing
character of existence
has not been fully
determined.

Third, the relationship
between invariance and
the ontological conditions
of continuity requires
further investigation.

Fourth, the status of
invariance across different
realms of being has not
yet been established.

Fifth, the relationship
between invariance and
the dynamics of being
remains open.

Sixth, the analysis has
not determined whether
there are conditions under
which invariance has a
more fundamental role
than the present relational
account can express.

These questions must not
be answered by importing
unexamined assumptions
about immutable substances,
permanent essences, or
fixed ontological foundations.

Their examination requires
the broader framework
that will emerge through
the subsequent development
of the book.

% ------------------------------------------------------------
% 10.11.1. TABLE: SYNTHESIS OF INVARIANCE
% ------------------------------------------------------------

% ============================================================
% TABLE 10 — SYNTHESIS OF INVARIANCE
% Book I: Ontology of Viability
% Part I: Identity and Manifestation
% Chapter 2: Invariance
% ============================================================

\begin{table}[htbp]
    \centering
    \caption{Conceptual Synthesis of Invariance}
    \label{tab:synthesis-of-invariance}

    \renewcommand{\arraystretch}{1.4}

    \begin{tabular}{|p{0.22\textwidth}|p{0.36\textwidth}|p{0.33\textwidth}|}
        \hline
        \textbf{Conceptual Relation} &
        \textbf{Established Distinction} &
        \textbf{Result for the Concept of Invariance} \\
        \hline

        Invariance and identity &
        Identity concerns the existent as identity-bearing; invariance concerns the remaining unchanged of a specified determination under specified transformations. &
        Invariance does not constitute identity or exhaust its meaning. \\
        \hline

        Invariance and difference &
        Difference concerns determinate distinctions, whereas invariance concerns the unchanged status of a particular respect within a specified domain. &
        Difference and invariance may coexist, provided that they concern distinct respects or determinations. \\
        \hline

        Invariance and change &
        Change concerns determinate alteration; invariance concerns the remaining unchanged of a specified determination under a specified transformation. &
        An existent may change in one respect while remaining invariant in another. \\
        \hline

        Invariance and continuity &
        Continuity concerns the continuation of an existent or process; invariance concerns the unchanged status of a particular determination. &
        Invariance does not, by itself, establish continuity, and continuity does not necessarily require a specified invariant feature. \\
        \hline

        Invariance and structure &
        Structure concerns an organized configuration of determinations and relations; structural invariance concerns the preservation of a specified structural determination or relation. &
        The invariance of a structural component does not establish the invariance of the whole structure. \\
        \hline

        Invariance and knowledge &
        Ontological invariance concerns the status of what remains unchanged; epistemic recognition concerns the ways in which a knower identifies and describes that invariance. &
        Knowledge does not constitute ontological invariance, and incomplete knowledge does not establish its absence. \\
        \hline

        Forms and domains of invariance &
        Quantitative, qualitative, relational, and structural forms are analytical distinctions; domains specify the transformations under consideration. &
        These forms are not assumed to be mutually exclusive or to constitute a final classification of ontological kinds. \\
        \hline

        Scope and limits &
        The relational account identifies a subject, a domain of transformations, and a respect in which something remains unchanged. &
        The account provides a working definition without settling the ultimate ontological status of invariance. \\
        \hline

    \end{tabular}
\end{table}

\vspace{0.5em}

\noindent
\textit{General synthesis:}
Invariance is provisionally understood in this chapter as a relation between an existent, or a determinate aspect of an existent, and a specified domain of changes or transformations, such that the existent or aspect remains unchanged in a particular respect under those transformations. This account distinguishes invariance from identity, difference, change, continuity, structure, and epistemic recognition without treating these concepts as wholly unrelated.

% ============================================================
% END OF TABLE 10
% ============================================================

% ------------------------------------------------------------
% 10.12. THE PLACE OF INVARIANCE
% IN THE DEVELOPMENT OF THE SYSTEM
% ------------------------------------------------------------

Invariance occupies a specific
but limited position within
the present development
of the ontology of viability.

It provides a means of examining
what remains unchanged in
particular respects when
an existent or one of its
determinations is considered
under specified transformations.

It therefore contributes
to the conceptual distinction
between identity, difference,
change, continuity, and
structure.

However, invariance does
not independently explain
the existence of an existent,
the origin of its identity,
or the conditions under
which it continues through
transformation.

These questions belong
to the broader ontological
framework and cannot be
settled through the concept
of invariance alone.

The present chapter has
therefore established a
disciplined conceptual
position rather than
a complete metaphysical
theory.

This limitation is deliberate.
The development of the
system requires that
conceptual distinctions
be established before
their deeper ontological
relations are determined.

The question of whether
invariance has a more
fundamental ontological
role may be reconsidered
when the relations among
the realms of being,
light, and the dynamics
of being have been
developed.

% ------------------------------------------------------------
% 10.13. TRANSITION TO THE NEXT CHAPTER
% ------------------------------------------------------------

The examination of invariance
has clarified one dimension
of the relationship between
an existent and the changes
or transformations under
which it is considered.

The next chapter turns
to the problem of unity
and plurality.

This transition is conceptually
necessary because the analysis
of invariance has repeatedly
encountered the coexistence
of different determinations,
relations, and structural
aspects within an existent.

The persistence of one
determination alongside
the alteration of another
raises questions concerning
how plurality is related
to the unity of what is
being examined.

The next chapter will
therefore investigate
the ontological problem
of unity and plurality
without presupposing
that unity must eliminate
difference or that plurality
must dissolve unity.

The distinction between
invariance and the concepts
examined in the present
chapter will remain relevant,
but the central inquiry
will now shift from what
remains unchanged under
transformation to the
relationship between
the one and the many.

% ============================================================
% END OF SECTION 10
% ============================================================

% ------------------------------------------------------------
% CHAPTER CONCLUSION
% ------------------------------------------------------------

\section*{Chapter Conclusion}
\addcontentsline{toc}{section}{Chapter Conclusion}

This chapter has approached invariance as a relational
problem within the ontology of viability. Its provisional
definition identifies invariance through the relation between
an existent or one of its determinate aspects and a specified
domain of change, with respect to which something remains
unchanged.

The analysis has distinguished invariance from identity,
difference, change, continuity, structure, and epistemic
recognition. These distinctions prevent the concept from
being reduced either to the persistence of an existent as
such or to the absolute absence of change.

The relational formulation also establishes an important
methodological restriction: a claim of invariance must
specify what is considered invariant, in relation to which
changes, and in what respect. Without these determinations,
the claim remains insufficiently defined.

At the same time, the relational definition does not settle
the ultimate ontological status of invariance. Whether
invariance is exclusively relational or may disclose a
deeper dimension of ontological determination remains an
open question for subsequent development and revision of
the philosophical system.

The chapter therefore establishes a conceptual framework
rather than a completed metaphysical explanation. Its
conclusions provide a basis for examining unity and plurality
in the next chapter, where the relations between unity,
multiplicity, distinction, and ontological determination
will be considered without presupposing that these problems
have already been resolved by the analysis of invariance.

% ============================================================
% END OF CHAPTER 02
% ============================================================



% ------------------------------------------------------------
% CHAPTER 03
% UNITY AND PLURALITY
% ------------------------------------------------------------

\input{2_MAIN_BODY/PART_01_IDENTITY_AND_MANIFESTATION/01_conceptual_exposition/CHAPTER_03_UNITY_AND_PLURALITY/chapter_03}


% ------------------------------------------------------------
% CHAPTER 04
% MANIFESTATION
% ------------------------------------------------------------

\input{2_MAIN_BODY/PART_01_IDENTITY_AND_MANIFESTATION/01_conceptual_exposition/CHAPTER_04_MANIFESTATION/chapter_04}


% ------------------------------------------------------------
% CHAPTER 05
% IDENTITY AND MANIFESTATION
% ------------------------------------------------------------

\input{2_MAIN_BODY/PART_01_IDENTITY_AND_MANIFESTATION/01_conceptual_exposition/CHAPTER_05_IDENTITY_AND_MANIFESTATION/chapter_05}


% ============================================================
% 02. FORMAL RECONSTRUCTION
% ============================================================

\section*{Formal Reconstruction}
\addcontentsline{toc}{section}{Formal Reconstruction}

The formal reconstruction examines the conceptual
commitments developed in the preceding chapters
and investigates their logical articulation.

Its purpose is to reconstruct the principles,
arguments, definitions, and formal relations
that can be justified on the basis of the
conceptual exposition.

Formal reconstruction must preserve the distinction
between established ontological commitments,
derived propositions, provisional hypotheses,
and unresolved questions.

Formal representation is not itself a substitute
for ontological justification. A formally coherent
expression does not become an established
ontological principle merely by virtue of its
formal consistency.

The chapter in this domain therefore examines
the logical structure of the account of identity
and manifestation while making its assumptions,
dependencies, and limits explicit.


% ------------------------------------------------------------
% CHAPTER 06
% PRINCIPLES, ARGUMENTS, AND FORMAL RECONSTRUCTION
% ------------------------------------------------------------

\input{2_MAIN_BODY/PART_01_IDENTITY_AND_MANIFESTATION/02_formal_reconstruction/CHAPTER_06_PRINCIPLES_ARGUMENTS_AND_FORMAL_RECONSTRUCTION/chapter_06}


% ============================================================
% 03. COMPARATIVE STUDIES
% ============================================================

\section*{Comparative Studies}
\addcontentsline{toc}{section}{Comparative Studies}

The comparative studies examine selected
ontological approaches in relation to the
conceptual problems developed in this part.

Comparison is intended to clarify the conceptual
position of the Viability Philosophy System,
identify points of convergence and divergence,
and examine the explanatory consequences of
different ontological commitments.

Comparative analysis must not assume that
different philosophical systems use identical
concepts merely because they employ similar
terminology.

The meaning, scope, and structural role of
each concept must be examined within its
respective theoretical context.

Nor should comparison be used to introduce
unexamined assumptions into the ontology of
viability. Its role is analytical and critical:
to clarify the developed framework and identify
questions that require further investigation.


% ------------------------------------------------------------
% CHAPTER 07
% COMPARATIVE ONTOLOGICAL STUDIES
% ------------------------------------------------------------

\input{2_MAIN_BODY/PART_01_IDENTITY_AND_MANIFESTATION/03_comparative_studies/CHAPTER_07_COMPARATIVE_ONTOLOGICAL_STUDIES/chapter_07}


% ============================================================
% PART CONCLUSION
% ============================================================

\section*{Conclusion}
\addcontentsline{toc}{section}{Conclusion}

Part 01 has established the conceptual and
analytical foundations for examining identity
and manifestation within the Ontology of Viability.

Its starting point is the ontological inseparability
of being and identity: to exist is to be
identity-bearing.

This commitment does not collapse the conceptual
distinction between being and identity. Instead,
it establishes the basis for examining identity
through its different conceptual dimensions
without treating it as an independent condition
added to existence.

The conceptual exposition develops the questions
of identity, invariance, unity and plurality,
manifestation, and the relation between identity
and manifestation.

The formal reconstruction examines the logical
structure of the resulting commitments and
distinguishes their formal articulation from
their ontological justification.

The comparative studies situate the developed
conceptual framework in relation to other
ontological approaches while preserving the
distinctive structure and terminology of the
Viability Philosophy System.

The conclusions of this part provide the
conceptual groundwork for the subsequent
examination of the realms of being.

However, the questions concerning the source
of identity, its relation to different
ontological levels, and the possibility of
higher-order identity manifesting through
an entity remain unresolved at this stage.

Their explanation requires the development
of the ontological architecture of the realms
of being and the investigation of light.
Only after these foundations have been
established can the dynamics of being
systematically address the constitution,
articulation, and manifestation of identity.

Part 01 therefore establishes the initial
conceptual foundation of identity and
manifestation without claiming to provide
a complete account of their ontological
constitution.


% ============================================================
% END OF PART 01
% ============================================================
